\documentclass[11pt]{article}
\usepackage[T1]{fontenc}
\usepackage{lmodern}
\usepackage[margin=1in]{geometry}
\usepackage{amsmath,amssymb,amsthm,mathtools,mathrsfs,amscd}
\usepackage{microtype}
\usepackage{enumitem}
\usepackage{booktabs}
\usepackage{xurl}
\usepackage[colorlinks=true,linkcolor=blue,citecolor=blue,urlcolor=blue]{hyperref}
\hypersetup{pdftitle={Unrestricted pro-modularity at the prime two},pdfauthor={OpenAI}}
\newtheorem{theorem}{Theorem}[section]
\newtheorem{proposition}[theorem]{Proposition}
\newtheorem{lemma}[theorem]{Lemma}
\newtheorem{corollary}[theorem]{Corollary}
\theoremstyle{definition}
\newtheorem{definition}[theorem]{Definition}
\theoremstyle{remark}
\newtheorem{remark}[theorem]{Remark}
\newcommand{\Q}{\mathbb Q}
\newcommand{\Z}{\mathbb Z}
\newcommand{\C}{\mathbb C}
\newcommand{\F}{\mathbb F}
\newcommand{\OO}{\mathcal O}
\newcommand{\kk}{k}
\newcommand{\eps}{\varepsilon}
\newcommand{\cts}{\mathrm{cts}}
\newcommand{\articleid}[1]{{\def\UrlBreaks{\do\-}\nolinkurl{#1}}}
\DeclareMathOperator{\GL}{GL}
\DeclareMathOperator{\SL}{SL}
\DeclareMathOperator{\PGL}{PGL}
\DeclareMathOperator{\Spec}{Spec}
\DeclareMathOperator{\Spf}{Spf}
\DeclareMathOperator{\Hom}{Hom}
\DeclareMathOperator{\End}{End}
\DeclareMathOperator{\tr}{tr}
\DeclareMathOperator{\ad}{ad}
\DeclareMathOperator{\HT}{HT}
\DeclareMathOperator{\Ind}{Ind}
\DeclareMathOperator{\rank}{rank}
\DeclareMathOperator{\edim}{edim}
\DeclareMathOperator{\Frac}{Frac}
\DeclareMathOperator{\Sym}{Sym}
\DeclareMathOperator{\WD}{WD}
\DeclareMathOperator{\Gal}{Gal}
\DeclareMathOperator{\Supp}{Supp}
\DeclareMathOperator{\Ann}{Ann}
\DeclareMathOperator{\Fil}{Fil}
\DeclareMathOperator{\gr}{gr}
\DeclareMathOperator{\coker}{coker}
\title{Unrestricted pro-modularity at the prime two}
\author{OpenAI}
\date{October 4, 2026}
\begin{document}
\maketitle
\begin{abstract}
Every continuous, odd, absolutely irreducible two-dimensional
$2$-adic representation of $G_{\Q}$ that is unramified outside finitely
many finite primes occurs in the full completed Hecke algebra at some
odd tame level. The level may contain auxiliary tame primes, and scalar
and reducible residual representations are included. This is a
completed-Hecke occurrence result, with no de Rham hypothesis; it does
not assert classical modularity.
\end{abstract}

\section{The full Hecke algebra and the theorem}
\label{sec:introduction}

Congruences between modular forms of different weights assemble their
Hecke eigenvalues into a single $p$-adic algebra. Its points need not
come from classical forms of one weight. The resulting modularity
question asks whether a Galois representation occurs in this larger
spectrum, even when no Hodge-theoretic condition is imposed at $p$.
We prove this occurrence statement for every odd, absolutely
irreducible two-dimensional representation over $\Q$ at $p=2$.

Fix an odd positive integer $N$. For $i\geq1$, let
$M_i(\Gamma_1(N))$ denote the space of all complex holomorphic modular
forms of weight $i$, including Eisenstein forms. Let
$T_{\leq k}^{(2)}(N)$ be the $\Z$-algebra acting on
$\bigoplus_{i=1}^{k}M_i(\Gamma_1(N))$
generated by $T_\ell$ and $\ell S_\ell$ for primes $\ell\nmid2N$.
On weight $i$, the operator $S_\ell$ is
$\ell^{i-2}\langle\ell\rangle$. Set
\begin{equation}\label{eq:full-hecke}
 \mathbb T_2(N)=
 \varprojlim_k\bigl(\Z_2\otimes_{\Z}T_{\leq k}^{(2)}(N)\bigr).
\end{equation}
Each finite-stage algebra has its $2$-adic topology, and
$\mathbb T_2(N)$ has the resulting inverse-limit topology.
We fix the Frobenius convention in which a classical eigenform has
characteristic polynomial
\begin{equation}\label{eq:frobenius}
 Z^2-T_\ell Z+\ell S_\ell
\end{equation}
at a good prime, as in \cite[Equation~(2.2)]{Dim}.

\begin{theorem}\label{thm:main}
Let $E/\Q_2$ be finite, with ring of integers $\OO_E$, and let
$r:G_{\Q}\longrightarrow\GL_2(E)$
be continuous and absolutely irreducible. Suppose that $r$ is
unramified outside finitely many finite primes and is odd:
$\det r(c)=-1$ for complex conjugation $c$.
Then there exist an odd positive integer $N$, divisible by every odd
prime where $r$ ramifies, and a continuous $\Z_2$-algebra homomorphism
\[
 \lambda:\mathbb T_2(N)\longrightarrow\OO_E
\]
such that, for every prime $\ell\nmid2N$,
\[
 \lambda(T_\ell)=\tr r(\operatorname{Frob}_\ell),\qquad
 \lambda(\ell S_\ell)=\det r(\operatorname{Frob}_\ell).
\]
\end{theorem}

We call this conclusion \emph{pro-modularity}. No condition on the
semisimplified reduction of $r$ is imposed. In particular, the theorem
includes scalar residual representations. The compact image of $r$
preserves an $\OO_E$-lattice, so its traces and determinants are integral.

\subsection*{Context and main ideas}

The Fontaine--Mazur conjecture relates Galois representations satisfying
geometric local conditions to algebraic geometry.  In its odd,
two-dimensional regular form over $\Q$, the expected conclusion is
classical cuspidal modularity up to a Tate twist
\cite{FM1995,Pan2022}.  Removing the local condition leads to a different
question: occurrence of the representation in a completed Hecke algebra.
Emerton records this expectation for the full varying-weight algebra in
\cite[Conjecture~2.12]{Emerton2011Families}.  That formulation prescribes
the allowed tame primes.  Theorem~\ref{thm:main} proves dyadic occurrence
for absolutely irreducible representations while allowing auxiliary
primes in the tame level.

Skinner and Wiles developed pro-modularity arguments through ordinary
families in the residually reducible setting
\cite[Introduction and Section~4.1]{SkinnerWiles1999}.
Emerton's local--global compatibility work separates pro-modularity from
the production of classical vectors
\cite[Section~1.2]{Emerton2011}; the pro-modularity theorem recorded there
assumes $p>2$, residual irreducibility, and local residual restrictions.
Pan extends the family method to the residually reducible nonordinary
setting at odd primes: ordinary loci supply modular points, and patching
at one-dimensional primes propagates Hecke support
\cite[Introduction and Theorem~4.1.7]{Pan2022}.
His classical modularity theorems retain regular de Rham hypotheses.
These precedents suggest placing an unrestricted representation in a
family and finding a useful geometric point elsewhere in that family.

The local representation theory used here comes from Colmez's
$p$-adic Langlands correspondence \cite{Colmez2010}, in the all-prime form
of Colmez, Dospinescu, and Pa\v{s}k\=unas
\cite[Theorem~1.1]{CDP2014}.  Pa\v{s}k\=unas and Tung supply the integral
block finiteness and Cayley--Hamilton descriptions needed at $2$,
including scalar residual blocks
\cite[Theorems~1.2--1.4 and Section~4.1]{PaskunasTung2021}.
Their results control local parameters but do not by themselves place a
global representation on Hecke support.

After a continuous cyclotomic twist, we place the original
representation in a fixed-determinant trace family $D$ of dimension
at least three. The curve and connectedness arguments of \cite{FM2}
put this family on Hecke support after solvable totally real base
change, as developed in Section~\ref{sec:family}. The main task is
then to find a regular de Rham point on $D$, from which Hecke
support can be transferred to the whole family over $\Q$.

Two constructions proved here produce that point. We choose the
base-change field to split completely at $2$, so that all dyadic
restrictions of the family have the same local Galois parameter.
Section~\ref{sec:diagonal} passes through the local block equivalences
to construct an object for one copy of $\GL_2(\Q_2)$ from an admissible
representation of the product of dyadic groups. Equality of the parameters
ensures that a failure of
admissibility would persist on a characteristic-two curve in every factor,
contradicting the admissibility of the product representation.
This comparison includes scalar residual blocks.
Section~\ref{sec:regular-point} then uses a chain of prime
specializations in $D$ to force maximal growth, and hence positive
Iwasawa rank, for the one-factor object. Polynomial approximation
detects a locally algebraic depth-zero supercuspidal type, yielding
a regular de Rham point $x$. The regular Fontaine--Mazur theorem
\cite[Theorem~1.1]{FM2} makes $x$ classical up to Tate twist, and
Section~\ref{sec:hecke} puts its classical packet and continuous
cyclotomic twists in the precise algebra \eqref{eq:full-hecke}.

The final passage uses the dimension theorem of
\cite[Theorem~1.1]{Dim} and the adjoint Selmer vanishing theorem of
Newton and Thorne \cite[Theorem~5.4]{NewtonThorne2023}.
At $x$, the global deformation problem with varying determinant has
tangent dimension at most three, while a dimension-four integral
Hecke component has local dimension three. Equality makes the
ambient local ring regular and forces its Hecke kernel to vanish
there. Since $D$ is a domain through $x$, the same kernel vanishes
on all of $D$. Section~\ref{sec:conclusion} carries out this comparison
and recovers the original coefficient point. The Hodge-theoretic
condition is needed only at the auxiliary point.

\subsection*{Coefficient and deformation conventions}

We allow finite extensions of coefficient fields during the proof.
Write $\OO$ for the current ring of integers, $\varpi$ for a
uniformizer, and $\kk$ for the residue field. A \emph{coefficient
point} of a complete local $\OO$-algebra is a continuous map to the
ring of integers of a finite coefficient extension. Its kernel is
a characteristic-zero prime; the corresponding integral quotient
has dimension one. A domain over $\OO$ is \emph{horizontal} if
$\varpi$ is nonzero in it.

At residue characteristic two, a pseudorepresentation means a
two-dimensional determinant law, with both trace and determinant
data, in the sense of \cite{Chenevier2014}. Thus no division by $2$
is implicit in an integral deformation ring. All global deformation
problems have a fixed finite ramification set containing $2$ and
allow ramification at infinity. Trace images and support containments
are closed images and set-theoretic containments; the domains used
for the family are reduced.

For the local categories and completed definite forms we use the
normalization of \cite[Section~2]{FM2}. If the fixed global
determinant is $\chi$, the central character is $\psi=\chi\eps$,
where $\eps$ is the cyclotomic character. The local block parameter
is the Galois parameter twisted by $\eps$, of determinant
$\psi\eps=\chi\eps^2$. Group duals have the contragredient action;
commuting coefficient and Hecke operators act by transposition.
The conversion to the Frobenius identity \eqref{eq:frobenius} is
made explicitly when passing to classical Hecke points.

\section{A trace family with potential automorphic support}
\label{sec:family}

We first place the given representation in a sufficiently large
fixed-determinant family.  We then put that entire family on completed
Hecke support after a totally real base change.  The representation
itself need not have an algebraic weight.  The distinction between its
residual representation and an auxiliary representation supplying the
determinant will be essential in the support argument.  For a matrix
representation, \emph{not virtually solvable} means that the identity
component of its algebraic monodromy group is not solvable.  This
property is unchanged on restriction to a finite index subgroup.

\subsection{Normalizing the determinant}

We will use the following twist property, proved in
Proposition~\ref{prop:hecke-twists}: at an odd tame level $N\geq5$, twisting
a Hecke coefficient point by any continuous character
$\eta:\Z_2^\times\to\OO'^\times$ gives another point of the same
Hecke algebra.  At a good prime $\ell$, this multiplies the trace by
$\eta(\eps(\operatorname{Frob}_\ell))$ and the determinant by its
square.  Thus it suffices to prove the theorem after any such twist
of $r$, and to undo the twist at the end.

\begin{lemma}\label{lem:normalization}
After a finite coefficient extension and a twist through $\eps$, one
may arrange
\[
 \det r=\chi=\delta\eps^w,
\]
where $\delta$ has finite order and is unramified at $2$.  The integer
$w$ may be chosen arbitrarily large in one parity class.  It may in
particular be chosen so that $\chi$ is the determinant of a regular
cuspidal modular representation that is not virtually solvable.
Oddness is preserved.  The character $\psi=\chi\eps$, viewed on
$\Z_2^\times$ by local reciprocity, is $z\mapsto z^{w+1}$.
\end{lemma}

\begin{proof}
Global class field theory separates the dyadic part of $\det r$ from
its odd conductor part:
\[
 \det r=\delta\,\alpha\circ\eps,
 \qquad \alpha:\Z_2^\times\longrightarrow\OO^\times.
\]
Here $\delta$ has finite order and odd conductor.  Indeed, at an odd
prime the pro-$\ell$ group of principal units has finite image in
$\OO^\times$, whose open subgroup is pro-$2$; only finitely many
primes occur.  The remaining unramified global character is trivial
by the class field theory of $\Q$.

Choose $w$ with $(-1)^w=\alpha(-1)$.  Write
$\Z_2^\times=\{\pm1\}\times5^{\Z_2}$.  The value $\alpha(5)$
reduces to one, since a pro-$2$ group has trivial image in
$\kk^\times$.  After finite coefficient extension choose
\[
 u^2=5^w\alpha(5)^{-1}.
\]
The element $u$ is integral and reduces to one.  Its powers $u^{2^n}$
tend to one, so $a\mapsto u^a$ defines a continuous character of
$\Z_2$.  Define $\beta(-1)=1$ and $\beta(5)=u$.  Then
$\alpha(z)\beta(z)^2=z^w$ for every $z\in\Z_2^\times$.
Replacing $r$ by $r\otimes(\beta\circ\eps)$ gives the asserted
determinant.  Since $\beta(-1)=1$, complex conjugation and oddness
are unchanged.  Local reciprocity identifies $\eps$ on dyadic units
with the identity character; $\delta$ is unramified there.

For completeness, the auxiliary modular representation can be chosen
without imposing the residual representation of $r$.  Oddness gives
$\delta(-1)=(-1)^{w+1}$.  Fix an odd level divisible by the conductor
of $\delta$, and enlarge it to a fine level.  The dimension formula
\cite[Section~III, Th\'eor\`eme~1]{CohenOesterle1977}
for cusp forms with this character grows linearly with the weight
$w+1$ along the indicated parity class.  The CM contribution at this
fixed level is bounded independently of the weight.  To see the
latter assertion, a CM eigenform comes from a Hecke character of an
imaginary quadratic field $K$; its level is
$|\operatorname{disc}K|\operatorname{Norm}(\mathfrak f)$, where
$\mathfrak f$ is the character conductor
\cite[Theorem~(3.4), Corollary~(3.5), and Theorem~(4.5)]{Ribet1977}.
Only finitely many pairs
$(K,\mathfrak f)$ occur at the fixed level.  For a fixed infinity
type the characters with conductor dividing a fixed ideal form, when
nonempty, a torsor under a finite ray class character group.  Oldform
multiplicities are bounded by the fixed level as well.  Thus a
sufficiently large $w$ gives a non-CM cuspidal eigenform of determinant
$\delta\eps^w$.  Its representation is regular and not virtually
solvable.  Indeed its restriction to every finite index subgroup is
semisimple; if its connected algebraic monodromy were solvable, its
unipotent radical would act trivially, so that connected group would
be a torus.  An open subgroup would then have abelian image, which
Ribet's CM criterion excludes
\cite[Propositions~(4.2) and~(4.4)]{Ribet1977}.
\end{proof}

We henceforth use the normalized $r$ and enlarge the finite
coefficient ring $\OO$ whenever necessary.  Every enlargement is
finite.  When a domain splits after coefficient extension, choose a
component through the specified coefficient point.  Finite flat
scalar extension makes its minimal prime contract to the old minimal
prime, so its dimension is unchanged.  Reduced support containments
also persist under this operation.  Thus the original coefficient
point and the dimension bounds below are retained.

\subsection{The dimension of the trace family}

\begin{lemma}\label{lem:large-family}
Fix a finite ramification set $S$ containing $2$, infinity, and the
ramified primes of the normalized $r$.  There is a reduced
horizontal framed deformation component $B$ of determinant $\chi$
through $r$.  Its completed trace image $D$ is a complete local
$\OO$-domain, contains the coefficient point of $r$, and satisfies
\[
 \dim B\ge6,\qquad \dim D\ge3.
\]
Moreover, $B$ has a coefficient point whose representation is not
virtually solvable.  Consequently the restriction of this family to
any finite extension of $\Q$ is generically absolutely irreducible.
\end{lemma}

\begin{proof}
Choose a stable lattice of $r$ and let $R^{\square,\chi}_{\bar r,S}$
be its unrestricted framed deformation ring with determinant $\chi$.
The residue here is the actual lattice reduction; no semisimplicity
or irreducibility is imposed on it.  These rings and their maps to
fixed-determinant pseudodeformation rings exist under Mazur's
finiteness condition, which holds for $G_{\Q,S}$
\cite[Proposition~3.1]{FM2}.

Complete at the characteristic-zero point of $r$.  Its real local
framed deformation ring is smooth of dimension two: an odd
involution has eigenvalues $1,-1$, and its conjugacy orbit has
stabilizer the two-dimensional diagonal torus.  For the global ring
over this real local ring, the relative relation estimate of
\cite[Proposition~3.6]{FM2}, with no imposed dyadic condition, is
\[
 g-r_{\mathrm{rel}}\ge
 3-h^0\bigl((\ad^0r)^\vee(1)\bigr)=3.
\]
The equality follows from absolute irreducibility.  In characteristic
zero the trace pairing identifies $\ad^0r$ with its dual; a nonzero
invariant in its cyclotomic twist would give a nonzero map
$r\to r\otimes\eps$.  Such a map is an isomorphism, whereas its
determinants would require $\eps^2=1$.  Every component of the
completed local ring therefore has dimension at least $2+3=5$.
The dimension comparison at a coefficient point in
\cite[Lemma~3.4]{FM2} adds one on returning to a horizontal integral
component.  Choose such a component through $r$ and give it the
reduced structure; its ring $B$ has dimension at least six.

Let $R^{\mathrm{ps},\chi}_{S}$ be the global pseudodeformation ring
of the semisimplified lattice reduction, and set
\[
 D=R^{\mathrm{ps},\chi}_{S}/
       \ker\bigl(R^{\mathrm{ps},\chi}_{S}\longrightarrow B\bigr).
\]
Write $T(g)\in D$ for its universal trace at $g$.  The ring $D$
is a horizontal complete local domain.  Its image in $B$ is
closed: complete local rings with finite residue field are compact,
and their continuous images in Hausdorff rings are closed.  The
generic representation on $B$ is absolutely irreducible, since the
irreducible locus contains $r$.  Forgetting its single frame loses
at most three dimensions, by \cite[Lemma~3.2]{FM2}.  Hence
$\dim D\ge\dim B-3\ge3$.

We justify the last assertion by bounding the trace loci of the
virtually solvable representations.  The reducible locus is closed
and proper on $D$, because it does not contain $r$.  An absolutely
irreducible virtually solvable two-dimensional representation is
induced from a character of a quadratic field, or has finite
projective image.  This follows by considering the connected
solvable algebraic subgroup: a noncentral torus has two eigenlines
permuted by the whole group, while a central connected subgroup gives
finite projective image.  A nontrivial connected unipotent subgroup
would have a unique invariant line and force reducibility.

Only finitely many quadratic fields $K$ can occur.  The corresponding
quadratic self-twist is unramified outside $S$, so this is
Hermite--Minkowski with bounded local degree.  For each $K$, the
maximal abelian pro-$2$ quotient of $G_{K,S}$ has rank at most two:
class field theory bounds its free part by the dyadic unit groups,
of total rank $[K:\Q]=2$; class groups and non-dyadic units contribute
only finite groups.  The invariant part under the quadratic
involution has rank at least one, from the cyclotomic extension.
If the inducing character is $\theta$, fixing the determinant fixes
$\theta\theta^s$ on $G_K$.  Thus at most one free character parameter
remains.  More explicitly, fixing the character on the image of
$1+s$ fixes a sublattice of rank at least one; the finite index in
its saturation contributes an integral, finite extension, not an
extra parameter.  The resulting complete character rings have
integral dimension at most two.  There are only finitely many
possible residual characters, since their unordered pair is
prescribed by $\bar r|_{G_K}$.

Their images in the trace space are closed of dimension at most two.
Indeed, each generating character value satisfies
\[
 X^2-T(h)X+\chi(h)=0\qquad(h\in G_K).
\]
Finitely many generators of the abelian character group, followed by
complete Nakayama, make each character ring finite over its trace
image.  This is the finite-character argument of
\cite[Proposition~3.1(3)]{FM2}.

For finite projective image, the classification of finite subgroups
of $\operatorname{PGL}_2$ leaves cyclic and dihedral groups, already
covered by reducibility and induction, and the three exceptional
groups $A_4,S_4,A_5$.  Their exponents divide $M=60$.
For any matrix with projective order dividing $M$, the ratio $\zeta$
of its eigenvalues satisfies $\zeta^M=1$, and
\[
 \frac{T(g)^2}{\chi(g)}=2+\zeta+\zeta^{-1}.
\]
Consequently $T(g)$ is a root of the monic polynomial
\[
 \prod_{\zeta^M=1}
  \bigl(X^2-\chi(g)(2+\zeta+\zeta^{-1})\bigr)\in\OO[X].
\]
Its coefficients belong to $\OO$ because the product over roots of
unity is Galois invariant and integral.  The fixed-determinant pseudo ring is
topologically generated by finitely many traces: traces generate it
by the determinant-of-a-sum identity, and finitely many suffice by
Nakayama on its cotangent space.  Imposing these monic equations on
such generators gives a finite $\OO$-algebra.  Thus the exceptional
finite-projective trace locus has integral dimension at most one.

The union of these finitely many closed loci cannot contain the
generic point of $D$, whose dimension is at least three.  Pull it
back to $B$ and use coefficient-point density
\cite[Lemma~3.3(2)]{FM2} to obtain a point outside it.  Restriction to
any finite index subgroup preserves the nonsolvable identity
component of its monodromy, so that point remains absolutely
irreducible.  The same is then true generically on the restricted
family and on every component containing that family.
\end{proof}

\subsection{Potential support with independently prescribed residue}

A closed trace locus over a totally real field is \emph{pro-modular}
if it lies in the spectrum of a residual factor of the completed
Hecke algebra of definite quaternionic forms, with the prescribed
central character and one tame level.  It is \emph{potentially
pro-modular} if this holds after a finite totally real solvable
extension, split completely above $2$.  All containments here use
the reduced closed trace images.  This is the support convention of
\cite[Section~2.3]{FM2}.

The potential pro-modularity theorem in \cite[Theorem~7.1]{FM2}
assumes a regular de Rham target.  Our family need not contain such
a target at this stage.  We therefore use its curve and connectedness
arguments with the determinant and residue specified separately.

\begin{proposition}\label{prop:potential-support}
Let $B$ be the family in Lemma~\ref{lem:large-family}, with the
normalized determinant of Lemma~\ref{lem:normalization}.  There is a
solvable totally real Galois extension $F/\Q$, of even degree and
completely split at $2$, such that the closed trace image of
$B|_{G_F}$ is pro-modular at one fixed tame level and central
character $\psi|_{G_F}$.  This assertion includes scalar and
reducible semisimple residual representations.
\end{proposition}

\begin{proof}
The determinant and the residual representation enter this argument
separately. Lemma~\ref{lem:normalization} supplies a determinant arising
from a regular modular representation; the modular seed below will
instead lift the residual representation of $B$.
We connect a characteristic-two curve on the seed's Hecke support
to one in an unrestricted locus containing the restricted family.
The constraints are used only to select and connect these curves.
Once the second curve has potential Hecke support,
localized propagation puts the entire unrestricted locus on that support.

We use the following precise propagation input.
The localized propagation theorem \cite[Theorem~5.1]{FM2} takes a
totally real solvable field of even degree, completely split at $2$,
a finite allowed set, a determinant of the form $\chi$ supplied by
a regular target, and, separately, a semisimple residual
representation of determinant $\bar\chi$.  If a characteristic-two
curve in this pseudodeformation space is potentially pro-modular,
has non-virtually-solvable generic representation, and has finite
local images at the allowed non-dyadic places, then every
irreducible closed locus containing that curve is potentially
pro-modular.

Choose a regular cuspidal modular lift of this semisimple residue,
using \cite[Lemma~7.2]{FM2}.  That lemma includes character sums and
scalar residue.  Enlarge the initial allowed set to contain its
tame primes and those of $B$, $\chi$, and the auxiliary regular
representation, before making any field choices.  Field preparation \cite[Proposition~6.1]{FM2} provides
abelian totally real $2$-extensions
\[
 \Q\subset A\subset F_0
\]
in which $2$ splits completely.  Put $d_A=[A:\Q]$ and
$d=[F_0:\Q]$.  The numbers $u_A,u_0$ of allowed non-dyadic primes
are bounded by a constant $U$ while both degrees grow.  Absolutely
irreducible residue remains absolutely irreducible.  In the
reducible case write $\bar r=\bar\chi_1\oplus\bar\chi_2$ and
$\alpha=\bar\chi_1/\bar\chi_2$.  Distinct characters remain
distinct, and restriction from $A$ to $F_0$ is injective on the
residual extension groups for $\alpha$ and $\alpha^{-1}$.

When the global residue is reducible and its dyadic characters coincide,
the connectedness argument will
keep track of chosen local invariant lines at selected dyadic places.
These places are the marks. Field preparation lets us make their
number large while keeping their proportion among all dyadic places
small.

There are no marks if the residue is absolutely irreducible or
$\alpha|_{G_{\Q_2}}\ne1$.  Otherwise field preparation specifies
one dyadic place $v_0$ of $A$, and we mark all its $m=[F_0:A]$
extensions to $F_0$.  Thus $m/d=1/d_A$; put $m=0$ in the unmarked
case.  We choose the degrees so that
\begin{equation}\label{eq:family-degrees}
 d>u_0+4,\qquad
 m>u_0+4\ \text{if marked},\qquad
 2d-2m-4-u_0>d.
\end{equation}
These are the degree conditions in \cite[Section~7.2]{FM2}.
Any further fixed lower bound on $d_A$ is compatible with them.

We now construct two large loci over $A$.  Solvable base change and
Jacquet--Langlands transfer the modular seed to definite forms.
Regularity keeps this base change cuspidal
\cite[Lemma~4.6]{FM2}.  Its central character and the prescribed
$\psi$ agree modulo $\varpi$.  At a sufficiently deep tame level,
the integral function description and exact reduction of completed
forms therefore give a nonzero residual summand with central
character precisely $\psi$; this is the seed construction in
\cite[Section~7.4]{FM2}.  A component $C_{\mathrm{seed}}$ of its
Hecke support has
\[
 \dim C_{\mathrm{seed}}\ge1+2d_A
\]
by \cite[Proposition~4.4]{FM2}.

Independently choose an unrestricted fixed-determinant framed
component containing the restricted family $B|_{G_A}$, by taking a
minimal prime below its defining prime.  Its generic representation
is absolutely irreducible by Lemma~\ref{lem:large-family}.
The relative presentation over the real local factors gives the
component dimension at least
$1+2d_A+3-h^0((\ad^0)^\vee(1))$; the invariant term is at most three.
Forgetting the frame loses at most three more dimensions.
Thus its trace image $C_{\mathrm{fam}}$, which contains the
restricted trace family, has
\[
 \dim C_{\mathrm{fam}}\ge2d_A-2
\]
by \cite[Proposition~3.6, Lemma~3.2, and Section~7.4]{FM2}.

In both loci impose $\varpi=0$ and the residual trace of one
Frobenius element at each allowed non-dyadic place.  In the marked
case impose also the constant residual local pseudorepresentation
at $v_0$.  The latter costs at most $\kappa$ equations, where
$\kappa$ is the number of generators of the closed-point ideal in
one fixed dyadic pseudo ring; it is independent of $d_A$.  Put
$\kappa=0$ when unmarked.  Taking
$d_A>U+\kappa+3$, the cut in $C_{\mathrm{fam}}$ has dimension at
least
\[
 2d_A-3-u_A-\kappa>d_A;
\]
the cut in $C_{\mathrm{seed}}$ has at least as large a lower bound.
Over the abelian field $A$, the characteristic-two reducible locus
has dimension at most one, the dihedral locus at most $d_A$, and
finite-projective-image curves are constant
\cite[Lemma~3.8]{FM2}.  Curve avoidance
\cite[Lemma~3.3(1)]{FM2} therefore gives a
non-virtually-solvable characteristic-two curve in each cut.
The Frobenius conditions make their allowed non-dyadic local images
finite by \cite[Lemma~3.9]{FM2}.  Denote these curves by
$C_s$ and $C_f$.

It remains to transfer support from $C_s$ to $C_f$.  Over $F_0$,
consider the characteristic-two framed representation scheme of
the fixed semisimple residue, adjoining an invariant projective
line at every mark.  A closed residual point is called
\emph{transverse} if none of these marked lines is globally
invariant; without marks this is an empty condition.
We check that both curves enter transverse residual charts.
Without marks, any stable lattice works.  With marks, the local
pseudo at $v_0$ is constant and scalar after a character twist.
The curve representation thus has a local invariant line, after a
finite extension of its field.  For a nonzero vector $v$ on that
line, over the normalized curve ring $J=\kk'[[t]]$ form
\[
 \mathcal V=\sum_{g\in G_A}J\rho(g)v.
\]
Compactness bounds this module in a lattice, and absolute
irreducibility makes it a full stable lattice.  The vector $v$ is
primitive: otherwise every translate would belong to
$t\mathcal V$, and Nakayama would give $\mathcal V=0$.
Its reduced orbit spans $\mathcal V/t\mathcal V$, so its reduction
belongs to no global invariant line.

This avoidance survives restriction to $F_0$.  A nonsplit extension
of distinct residual characters remains nonsplit by the restriction
injectivity, while a split representation retains its two distinct
character lines.  For globally scalar semisimplification, the orbit
condition rules out a split scalar reduction; its nonzero additive
extension remains nonzero by the same injectivity.  These statements
hold after any residual field extension.  Finally, $F_0/A$ is
Galois, so transporting the local line to all conjugate marks
preserves avoidance of every global invariant line.  Both curves
therefore pass through transverse residual points, including in
the scalar residual case.

We recall what the two connectedness inputs prove, to make the
support transfer explicit.  At a transverse point,
\cite[Lemma~7.3]{FM2} gives components of the completed chart of
dimension at least $2d-m$, with successive intersections in a
connecting chain of dimension at least $2d-m-1$.  The reducible
locus has strictly smaller dimension.  For an intersection
component of maximal dimension, forgetting its frame and marked
lines loses at most $3+m$ dimensions, and forgetting the characters
on those lines is finite.  Its trace image therefore has dimension
at least $2d-2m-4$.  The $u_0$ Frobenius trace conditions leave
dimension greater than $d$ by \eqref{eq:family-degrees}, so the same
bad-locus avoidance supplies a non-virtually-solvable curve with
finite bad local images.  Localized propagation carries support
across each such intersection.  This is
\cite[Lemma~7.4]{FM2}.

The transverse residual fiber is geometrically connected by
\cite[Lemma~7.5]{FM2}.  In particular this includes scalar residue:
after fixing the unique global invariant line, the nonzero additive
extension classes and the transverse marked lines give a connected
parameter space.  The finite graph of components meeting this fiber
therefore joins the seed chart to the family chart.  Starting from
$C_s$, which is pro-modular, the preceding chart propagation makes
$C_f|_{G_{F_0}}$ potentially pro-modular.  Hence $C_f$ over $A$ is
potentially pro-modular.  Apply localized propagation now over $A$
to the \emph{unrestricted} locus $C_{\mathrm{fam}}$ containing
$C_f$.  The entire locus, and therefore the prescribed family,
acquires Hecke support.  Marks have served only to connect curves
in characteristic two; they impose no condition on $B$.

Only finitely many components, charts, and propagation steps occur.
Taking the common Galois compositum of the resulting fields and a
common deeper tame level, as in \cite[Lemma~4.6 and Section~7.4]{FM2},
gives one extension $F/\Q$ with all the stated properties.
\end{proof}

\subsection{The diagonal local parameter and completed forms}

Write $D$ for the trace domain of Lemma~\ref{lem:large-family} and
choose $F$ from Proposition~\ref{prop:potential-support}.  Let
$\mathbb T$ be the residual completed Hecke algebra over $F$ just
obtained.  Support containment gives a continuous map
$\mathbb T\to D$: its image is the closed algebra of traces
restricted to $G_F$.  Indeed the Hecke kernel vanishes in this
reduced domain, by the definition of closed support.

Let $R$ be the local pseudodeformation ring of the residual dyadic
parameter, with determinant $\chi|_{G_{\Q_2}}\eps^2$.  This is
\emph{block normalization}: we twist each Galois parameter by
$\eps$, and the central character is
$\psi=\chi\eps$.  Because $2$ splits completely in $F$, each of
the $b=[F:\Q]$ dyadic restrictions of the global family has this
same parameter as a $D$-valued determinant.  The identifications use conjugate decomposition
groups; traces are unchanged by the conjugations.

For the corresponding compact dual block of
$G=\GL_2(\Q_2)$ with this central character, let $P$ be the direct
sum of one projective envelope of the dual of each simple object,
and put $\mathcal E=\End(P)$.  Coefficients are enlarged so that
the block is absolutely split.  The projective-generator functor
$\Hom(P,-)$ takes values in right $\mathcal E$-modules; its inverse
is completed tensor product with $P$.  We use the exact compact
block equivalence here, and recall its finer local properties in
Section~\ref{sec:diagonal}.

The next proposition puts the trace family and these local actions on
one coefficient module. Its finiteness over $D$ and $R$ will let us
compare specialization with admissibility.

\begin{proposition}\label{prop:family}
For the normalized representation $r$, the preceding constructions
give a horizontal complete local domain $D$ of dimension at least
three containing its coefficient point, and continuous maps
\[
 \mathbb T\longrightarrow D,\qquad R\longrightarrow D
\]
that make $D$ finite over both source rings.  There is a right
$\mathcal E^{\widehat\otimes b}$-module $L_0$, finite and faithful
over $\mathbb T$, whose inverse under the product block equivalence
is the residual summand of compact definite forms.  The module
\[
 L=\bigl(L_0\otimes_{\mathbb T}D\bigr)
       \big/\bigl(\text{$D$-torsion}\bigr)
\]
is nonzero, finite and torsion-free over $D$.  All $b$ central
$R$-actions on $L$ are the same map $R\to D$.  Its inverse under
the product block equivalence is finite over the completed group
algebra of a sufficiently small product of determinant-one compact
subgroups.
\end{proposition}

\begin{proof}
Only finiteness and the module assertions remain to be proved.
Let $D_F$ be the closed restricted trace image in $D$.  If
$n=[F:\Q]$, then $g^n\in G_F$ for every $g\in G_\Q$.
The trace-of-powers recurrence
\[
 P_0(X,a)=2,\quad P_1(X,a)=X,\quad
 P_j(X,a)=XP_{j-1}(X,a)-aP_{j-2}(X,a)
\]
gives a monic polynomial $P_n(X,a)$ of degree $n$, with
\[
 P_n(T(g),\chi(g))=T(g^n)\in D_F.
\]
Choose finitely many traces topologically generating $D$ over
$\OO$.  Each is integral over $D_F$, so the subalgebra they generate
is finite over $D_F$.  It is compact and hence closed in $D$;
since it also contains a dense subalgebra, it equals $D$.
Thus $D$ is finite over $D_F$, and therefore over $\mathbb T$.

The local-global block comparison \cite[Proposition~4.2]{FM2}
makes $\mathbb T$ finite over
$R^{\widehat\otimes b}$, with its actual Galois restriction action
in block normalization.  On $D$ this action factors through
\[
 R^{\widehat\otimes b}\longrightarrow R,
 \qquad a_1\otimes\cdots\otimes a_b\longmapsto a_1\cdots a_b,
\]
because all dyadic parameters agree.  Transitivity of finiteness
then proves that $D$ is finite over the single ring $R$.

Let $M$ be the compact definite-forms summand.  The same proposition
gives
\[
 L_0=\Hom_{G^b,\mathrm{cts}}(P^{\widehat\otimes b},M),
\]
finite and faithful over $\mathbb T$, with the asserted compatible
product action.  The compact module $M$ is admissible, equivalently
finite over a sufficiently deep determinant-one completed group
algebra, by \cite[Proposition~4.1]{FM2}.

At the generic point of $D$, the base change of $L_0$ is nonzero.
Indeed, if $\mathfrak q$ is its contraction to $\mathbb T$,
faithfulness and finiteness give
$\operatorname{Supp}_{\mathbb T}L_0=\Spec\mathbb T$; Nakayama
then gives $L_0\otimes_{\mathbb T}\kappa(\mathfrak q)\ne0$.
Extending this residue field to $\Frac(D)$ preserves nonvanishing.
Quotienting by $D$-torsion consequently leaves a nonzero finite
$D$-module $L$.  The torsion submodule is stable under every
$\mathcal E$-action because those actions commute with $D$.
Since $D$ is horizontal, $L$ is also $\OO$-torsion-free.

Finally, finiteness of $D$ over $\mathbb T$ makes
$L_0\otimes_{\mathbb T}D$, and then $L$, a quotient of a finite
direct sum of copies of $L_0$ as product-block modules.  These are
maps of compact modules with closed images.  Exactness of the
product block equivalence makes the inverse object of $L$ a
quotient of the corresponding finite sum of copies of $M$.
It is therefore admissible as asserted.  This supplies product
admissibility; passing to a single factor is the separate argument
of the next section.
\end{proof}

\section{Admissibility from diagonal local parameters}\label{sec:diagonal}

Proposition~\ref{prop:family} gives an admissible compact object for
\(G^b\), where \(G=\GL_2(\Q_2)\) and \(b=[F:\Q]\). We need an
admissible object for one copy of \(G\), still carrying the whole trace
family \(D\). The equality of the \(b\) local Galois parameters is the
additional information that makes this possible.

Retain the local block generator \(P\), its endomorphism algebra
\(\mathcal E=\End(P)\), and its fixed-determinant pseudodeformation ring
\(R\) from Proposition~\ref{prop:family}. Thus the central
character is \(\psi=\chi\eps\), the Galois determinant in block
normalization is \(\psi\eps\), and the compact category uses the dual
central action. Write
\[
 \mathcal E^{(b)}=\widehat\bigotimes_{i=1}^b\mathcal E,
 \qquad P^{(b)}=\widehat\bigotimes_{i=1}^bP.
\]
All completed tensor products in this section are in the category of
pseudocompact modules. Coinvariants mean quotients by the closed
submodule generated by the indicated relations.

\begin{proposition}[One-factor admissibility]\label{prop:diagonal}
Let \(D\) be a complete noetherian local \(\OO\)-domain in which
\(\varpi\ne0\), and suppose that \(R\to D\) is a finite continuous
local map. Let \(L\ne0\)
be a finite torsion-free \(D\)-module with a commuting continuous right
\(\mathcal E^{(b)}\)-action. Suppose that every copy of \(R\) acts
through the same map \(R\to D\), and that
\[
 M=L\widehat\otimes_{\mathcal E^{(b)}}P^{(b)}
\]
is admissible in the compact category. Then, using the first
\(\mathcal E\)-action,
\[
 V=L\widehat\otimes_{\mathcal E}P
\]
is nonzero, \(\OO\)-torsion-free, and finite over
\(\mathscr A=\OO[[K]]\), for every sufficiently small uniform open
subgroup \(K\subset\SL_2(\Z_2)\). It has a continuous commuting
\(D\)-action.
\end{proposition}

The construction takes place on the algebra side of the block equivalence:
\[
\begin{CD}
 L_{\mathcal E^{(b)}} @>{-\widehat\otimes_{\mathcal E^{(b)}}P^{(b)}}>> M\\
 @V{\text{forget factors }2,\ldots,b}VV @.\\
 L_{\mathcal E} @>{-\widehat\otimes_{\mathcal E}P}>> V.
\end{CD}
\]
The subscripts indicate which algebra acts on \(L\). The lower row
reconstructs the one-factor object from that algebra module.

Here admissibility in the compact category means finite generation over
the corresponding completed compact group algebra. Although \(K\) is
not open in \(G\), its product with the scalar units is open in
\(\GL_2(\Z_2)\); the fixed central character therefore makes this the
usual admissibility condition. We prove the proposition by comparing
one-factor and product coinvariants. The essential local assertion will
be that, along a nonconstant characteristic-two curve, all simple block
modules give the same answer to the question whether their
\(K\)-coinvariants vanish.

\subsection{The local block input}

We record the precise consequences of Pa\v{s}k\=unas--Tung that enter the
argument. The residue field has been enlarged so that the block and its
simple objects are absolutely split. Let \(e\in\mathcal E\) be the sum
of the idempotents corresponding to those simple objects on which
\(\SL_2(\Q_2)\) does not act trivially.

\begin{theorem}[Pa\v{s}k\=unas--Tung]\label{input:local-blocks}
The following statements hold for every such block at \(p=2\),
including the block of a scalar residual Galois representation.
\begin{enumerate}
\item The functor \(\Hom(P,-)\) is an exact equivalence from the compact
block to right pseudocompact \(\mathcal E\)-modules, with inverse
\(-\widehat\otimes_{\mathcal E}P\). The algebra \(\mathcal E\) is finite
over \(R\).
\item Let \(\mathrm{CH}_R\) be the universal Cayley--Hamilton algebra
for the local pseudorepresentation, and let
\(\mathrm{CH}_R^{\mathrm{tf}}\) be its maximal
\(\OO\)-torsion-free quotient. There is an \(R\)-compatible
identification
\[
 \mathrm{CH}_R^{\mathrm{tf}}\simeq
 \begin{cases}
  (e\mathcal E e)^{\mathrm{op}},&\text{in a nonsupersingular block},\\
  M_2((e\mathcal E e)^{\mathrm{op}}),&\text{in a supersingular block}.
 \end{cases}
\]
The corner equivalence is induced by the Montr\'eal functor on the
quotient by the objects with trivial \(\SL_2(\Q_2)\)-action.
\item An \(\mathcal E\)-module finite and torsion-free over \(\OO\) gives,
after duality and inversion of \(2\), an admissible Banach representation
of finite length. At an absolutely irreducible Galois coefficient
parameter there is a unique absolutely irreducible Banach object. At a
sum of characters \(\gamma_1+\gamma_2\), the simple objects have the
following description. If the characters coincide, the unique simple is
\(\Ind_B^G(1\otimes\eps^{-1})^{\cts}\otimes\gamma_1\circ\det\).
If their sum is \(\gamma+\gamma\eps\), the simples are
\[
 \{1,\widehat{\mathrm{St}},
       \Ind_B^G(\eps\otimes\eps^{-1})^{\cts}\}
       \otimes\gamma\circ\det,
\]
where \(\widehat{\mathrm{St}}\) is continuous Steinberg. If
\(\gamma_1\gamma_2^{-1}\notin\{1,\eps,\eps^{-1}\}\), they are exactly
\[
 \Ind_B^G(\gamma_1\otimes\gamma_2\eps^{-1})^{\cts},
 \qquad
 \Ind_B^G(\gamma_2\otimes\gamma_1\eps^{-1})^{\cts}.
\]
Here \(B\) is the upper Borel, induction is unnormalized, and Galois
characters are viewed as characters of \(\Q_2^\times\) by reciprocity.
\end{enumerate}
\end{theorem}

The first assertion is \cite[Section~4.1 and Theorems~1.2--1.3]
{PaskunasTung2021}; the second is
\cite[Propositions~4.10 and~4.18, Theorem~6.13]{PaskunasTung2021};
the Banach assertions are \cite[Section~4.5, Section~6.2,
Equation~(20), Proposition~6.9, Corollary~6.10, and
Proposition~6.11]{PaskunasTung2021}.
At \(2\), the center comparison gives an isomorphism after inverting
\(2\) and permits an integral cokernel killed by \(2\)
\cite[Theorem~1.4]{PaskunasTung2021}. We use the finite map from \(R\)
and the stated integral Cayley--Hamilton comparison, without identifying
\(R\) with the integral center.

The compact equivalence commutes with coefficient quotients, finite
presentations, and inverse limits. Consequently it preserves injections
and coefficient torsion-freeness, and commutes with the completed
coefficient changes below. The corresponding product equivalence is
explained in \cite[Proposition~4.2]{FM2}.

\subsection{Coinvariants and the noncharacter corner}

Fix a sufficiently small uniform \(K\subset\SL_2(\Z_2)\), and put
\[
 C_K=P/(\varpi P+(K-1)P).
\]
This is a left \(\mathcal E\)-module. Its support over \(R\) will control
the admissibility test.

\begin{lemma}\label{lem:coinvariant-finite}
The module \(C_K\) is finite over \(R\). For a finite pseudocompact
right \(\mathcal E\)-module \(W\),
\[
 (W\widehat\otimes_{\mathcal E}P)/(\varpi,K-1)
       =W\otimes_{\mathcal E}C_K.
\]
The analogous formula holds for product coinvariants. If a commutative
complete coefficient algebra $A$, finite over $R$ (over
$R^{\widehat\otimes b}$ in the product case), acts on $W$ compatibly
and commutes with its algebra action, these formulas commute with
base change $A\to H$ to a field.
\end{lemma}
\begin{proof}
The object \(P/\mathfrak m_RP\) corresponds under the compact
equivalence to \(\mathcal E/\mathfrak m_R\mathcal E\), so has finite
length. Every smooth simple in the block is admissible. Its compact
dual therefore has finite-dimensional \(K\)-coinvariants, with the
central character understood. Right exactness of coinvariants now
makes \(C_K/\mathfrak m_RC_K\) finite-dimensional. The
\(\mathfrak m_R\)-action is topologically nilpotent: on the algebra
side this follows from the finite \(R\)-module \(\mathcal E\), and the
equivalence transports the inverse limit of its finite-length
quotients. Compact Nakayama proves that \(C_K\) is finite over \(R\).

Completed tensor products commute with the cokernels defining
coinvariants. After these cokernels, all modules involved are finite
over their complete noetherian coefficient rings, so ordinary and
completed tensor products agree. For the product formula, apply this
argument in each factor. Finite presentations over \(R\), or over its
completed tensor powers, then show that the formula commutes with the
asserted field base changes.
\end{proof}

A characteristic-two curve in \(\Spec R\) means the spectrum of a
one-dimensional complete local domain quotient of \(R/\varpi\). Its
generic point is nonconstant if its map to \(\Spec R\) is not the
closed residual point.

\begin{lemma}\label{lem:full-corner}
At the generic point of a nonconstant characteristic-two curve, the
idempotent \(e\) is full: if \(H\) is a finite extension of the curve's
fraction field and \(\mathcal E_H=\mathcal E\otimes_RH\), then
\(\mathcal E_He\mathcal E_H=\mathcal E_H\).
\end{lemma}
\begin{proof}
We show that \(\mathcal E/\mathcal E e\mathcal E\) is finite over
\(\OO\). Under the compact equivalence its modules have only the
\(\SL_2(\Q_2)\)-trivial simple constituents. This is precisely the
subcategory of objects with trivial \(\SL_2(\Q_2)\)-action: that
subcategory is thick, and the assertion passes to compact limits
\cite[Section~4.2]{PaskunasTung2021}.

Such an action factors through \(\det:G\to\Q_2^\times\). On squares,
the character is prescribed by the central action. Since
\(\Q_2^\times/(\Q_2^\times)^2\) is finite, the corresponding completed
group algebra with these scalar relations is finite over \(\OO\).
This remains true if the scalar relations force coefficient torsion.
The regular module of \(\mathcal E/\mathcal E e\mathcal E\) corresponds
to an object of this subcategory with finite cosocle. Compact Nakayama
makes it finite over that finite \(\OO\)-algebra; its endomorphism
algebra, and hence \(\mathcal E/\mathcal E e\mathcal E\), is finite over
\(\OO\). Its special-fiber support over \(R\) is therefore contained in the
closed point. It vanishes at the curve's generic point, as required.
\end{proof}

\subsection{The characteristic-two curve test}

Let $H$ be a finite extension of the fraction field of a nonconstant
characteristic-two curve, chosen to split the parameter and the algebra
$\mathcal E_H=\mathcal E\otimes_R H$. We will show that a nonzero
$C_{K,H}=C_K\otimes_R H$ survives tensoring with every simple right
$\mathcal E_H$-module. The main case is a parameter
with two distinct characters. Its two principal-series families have
the same coinvariant vanishing: exchanging their inducing characters
inverts their values on determinant-one stabilizers. For a unit $a$,
the elements $a-1$ and $a^{-1}-1$ generate the same ideal.
To apply this calculation to all block simples, we construct the
families integrally and track their two character labels through
specialization.

\begin{lemma}\label{lem:principal-families}
Suppose a nonconstant characteristic-two curve has generic local
parameter \(\gamma_1+\gamma_2\), with \(\gamma_1\ne\gamma_2\). After a
finite extension, let \(J=\kk'[[t]]\) be its normalization, finite
over the curve ring, and
\(H=\Frac(J)\). The compact dual principal-series families over \(J\)
with inducing orders
\[
 (\gamma_1,\gamma_2\eps^{-1}),\qquad
 (\gamma_2,\gamma_1\eps^{-1})
\]
correspond to finite free \(J\)-modules \(W_1,W_2\) for
\(\mathcal E\). Their central \(R\)-action is the given curve parameter.
Their generic fibers are the two distinct simple right
\(\mathcal E_H\)-modules, and
\[
 (W_1\otimes_JH)\otimes_{\mathcal E_H}(C_K\otimes_RH)\ne0
 \quad\Longleftrightarrow\quad
 (W_2\otimes_JH)\otimes_{\mathcal E_H}(C_K\otimes_RH)\ne0.
\]
\end{lemma}
\begin{proof}
The characters take values in \(J^\times\): their values are integral
by the characteristic polynomial, and their product is a unit. They
are continuous. Indeed, choose an element on which the two characters
differ; traces against that element recover their values continuously
in \(H\). Their reductions are the residual constituents.

We first construct the families in mixed characteristic. Enlarge the
integer coefficients to \(\OO'\), and set
\(J_0=\OO'[[X,Y]]\). The pro-\(2\) completion of
\(\Q_2^\times\) has two free generators, represented by \(2\) and
\(5\), and torsion generated by \(-1\). Let \(\Gamma_1\) be the
universal lift of the first residual character on the two free
generators, fixing either lift of its sign on \(-1\), and put
\(\Gamma_2=\psi\eps\Gamma_1^{-1}\). There is a continuous map
\(J_0\to J\) specializing these characters to \(\gamma_1,\gamma_2\).

For each order, form smooth induction with coefficients in the
discrete dual \(J_0^\vee\), and take its compact dual \(N_i\).
This defines the family even when the inducing characters become
smooth only modulo powers of the coefficient maximal ideal.
Local sections on \(B\backslash G=\mathbf P^1(\Q_2)\) identify its
compact model with a pro-free \(J_0\)-module. In particular the model
is flat and commutes with coefficient changes. The residual smooth
induction has finite length in the specified block. If the residual
characters differ it is irreducible; if they coincide it has the
character and Steinberg constituents. In either case it has precisely
one noncharacter constituent, with multiplicity one
\cite[Section~4.1]{PaskunasTung2021}. Filtering the finite coefficient
quotients shows that \(N_i\) belongs to the compact block.

Put \(\widetilde W_i=\Hom(P,N_i)\). Compact Nakayama and residual
finite length show that this is a finite \(J_0\)-module. Exactness of
the equivalence makes it flat: for every finitely generated ideal
\(I\subset J_0\), apply the equivalence to the injective map
\(I\otimes_{J_0}N_i\to N_i\), using finite presentations to commute
these tensors with the equivalence. Thus \(\widetilde W_i\) is finite
free. The same argument proves compatibility with coefficient changes.
The residual multiplicity calculation gives
\[
                    \rank_{J_0}(\widetilde W_i e)=1.
\]

At characteristic-zero coefficient points outside the proper loci
\(\Gamma_1\Gamma_2^{-1}\in\{1,\eps,\eps^{-1}\}\), Schikhof duality
identifies these objects with the two continuous principal series in
Theorem~\ref{input:local-blocks}. Those points are Zariski dense in
\(\Spec J_0\). Hence their central \(R\)-action is everywhere the
pseudorepresentation \(\Gamma_1+\Gamma_2\): equality may be tested on
matrices acting on the finite free modules.

We also need to keep the two character labels distinct upon
specialization. The residual block here is nonsupersingular, so
Theorem~\ref{input:local-blocks} identifies the opposite corner with
\(\mathrm{CH}_R^{\mathrm{tf}}\), without a matrix factor.
Its Galois action on each free rank-one module
\(\widetilde W_i e\) gives a character. At the dense coefficient
points the two principal series are distinct simple modules, their
nonzero corner modules are distinct, and their corner labels are
therefore \(\Gamma_1\) and \(\Gamma_2\). More explicitly, if
\(\lambda_i\) is the corner character, set
\[
 I_{ij}=(\lambda_i(g)-\Gamma_j(g):g\in G_{\Q_2})\subset J_0.
\]
The dense coefficient points lie in
\(V(I_{i1})\cup V(I_{i2})=V(I_{i1}I_{i2})\), so
\(I_{i1}I_{i2}=0\). Since \(J_0\) is a domain, one of the two ideals
vanishes. The two labels are opposite because the generic
specializations are distinct. Thus the labels on the entire family
are \(\Gamma_1,\Gamma_2\), up to one fixed permutation, and these
identities persist under \(J_0\to J\).

Set \(W_i=\widetilde W_i\otimes_{J_0}J\). By
Lemma~\ref{lem:full-corner}, \(e\) is full over \(H\); its rank-one
corner modules therefore make \(W_i\otimes_JH\) simple and distinct.
The specialized torsion-free Cayley--Hamilton algebra is a quotient
of the Cayley--Hamilton algebra of the specialized parameter. The
field structure theorem for a split multiplicity-free determinant
allows at most the two character simples
\cite[Theorems~2.12 and~2.22]{Chenevier2014}. Thus our two modules
account for all simples, without assuming that forming the torsion-free quotient
commutes with reduction modulo \(2\).

Finally calculate \(K\)-coinvariants directly in the compact
induction model. There are finitely many double cosets
\(B\backslash G/K\). For a representative \(g\), its contribution
is \(J\) modulo the ideal generated by the inducing character minus
one on \(B\cap gKg^{-1}\). This follows equally by dualizing the
invariant functions with values in \(J^\vee\). Every stabilizer has
determinant one. Since \(\eps=1\) over \(J\), the two inducing
characters on that stabilizer are inverse. They generate the same
ideal, because \(a^{-1}-1=-a^{-1}(a-1)\) for a unit \(a\).
The two coinvariant modules consequently vanish simultaneously after
tensoring with \(H\). Lemma~\ref{lem:coinvariant-finite} identifies
these modules with the two displayed tensor products.
\end{proof}


\begin{lemma}[Uniform survival on curves]\label{lem:curve-support}
Let \(H\) be a finite extension of the fraction field of a nonconstant
characteristic-two curve in \(\Spec R\), large enough to split the
parameter and the finite algebra \(\mathcal E_H\). If
\(C_{K,H}=C_K\otimes_RH\ne0\), then
\[
                  S\otimes_{\mathcal E_H}C_{K,H}\ne0
\]
for every simple right \(\mathcal E_H\)-module \(S\).
\end{lemma}
\begin{proof}
The corner is full by Lemma~\ref{lem:full-corner}. The
Cayley--Hamilton comparison and its field structure therefore give
at most the simples labelled by the absolute constituents of the
parameter. Repeated characters cannot occur on this curve. Indeed
the fixed determinant has constant finite values in characteristic
two; if it equals \(\gamma^2\), injectivity of squaring in a field
makes \(\gamma\), and hence the whole parameter, constant.

If the parameter is absolutely irreducible, the split algebra has only
one simple module. A nonzero finite left module has a nonzero
semisimple head, so tensoring it with that simple right module is
nonzero. If the parameter has two distinct characters,
Lemma~\ref{lem:principal-families} constructs both simple
modules and proves that their tensor tests have identical vanishing.
At least one test is nonzero: otherwise the semisimple head of \(C_{K,H}\) would be
zero. Both tests are therefore nonzero.
\end{proof}

\subsection{From product coinvariants to one-factor admissibility}

\begin{proof}[Proof of Proposition~\ref{prop:diagonal}]
The module \(L\) is finite over \(R\), hence over the first
\(\mathcal E\). The exact compact equivalence makes \(V\) nonzero and
\(\OO\)-torsion-free. It also transports the continuous commuting
\(D\)-action. By Lemma~\ref{lem:coinvariant-finite},
\[
 Q:=V/(\varpi,K-1)=L\otimes_{\mathcal E}C_K
\]
is finite over \(D/\varpi\). It suffices, by compact Nakayama over
\(\mathscr A\), to show that \(Q\) is finite-dimensional over \(\kk\).

Suppose otherwise. Its closed support contains a one-dimensional
complete local domain quotient of \(D/\varpi\). At the generic point
of this curve, \(Q\) has nonzero fiber. Since \(D\) is finite over
\(R\), its image in \(\Spec R\) is also a nonconstant curve. Extend
its fraction field to a field \(H\) as in
Lemma~\ref{lem:curve-support}. With subscripts denoting fibers, we have
\[
                 L_H\otimes_{\mathcal E_H}C_{K,H}\ne0.
\]
This tensor product still has the commuting actions of the remaining
\(b-1\) copies of \(\mathcal E_H\).

For any nonzero finite right \(\mathcal E_H\)-module \(U\), choose a
simple quotient \(S\). Right exactness of tensor products and
Lemma~\ref{lem:curve-support} give
\(U\otimes_{\mathcal E_H}C_{K,H}\ne0\). Apply this observation
successively to the remaining factors. The central parameters are
equal, so the same algebra and the same module \(C_{K,H}\) occur each
time. We obtain
\[
 L_H\otimes_{\mathcal E_H^{\otimes_H b}}C_{K,H}^{\otimes_H b}\ne0.
\]
By the product formula in Lemma~\ref{lem:coinvariant-finite}, this is
the fiber of \(M/(\varpi,K^b-1)\) on the chosen curve. But product
admissibility makes that module finite-dimensional over \(\kk\).
Its continuous \(D\)-action is supported at the closed point, so its
fiber on a nonconstant curve is zero. This contradiction proves that
\(Q\) is finite-dimensional and hence that \(V\) is finite over
\(\mathscr A\).
\end{proof}

\section{A regular point in the trace family}\label{sec:regular-point}

The admissible module of Proposition~\ref{prop:diagonal} still need not
have an algebraic vector.  We first show that the size of its commuting
trace family forces positive Iwasawa rank.  That rank supplies algebraic
vectors with a prescribed smooth type.  A coefficient eigensystem among
these vectors will give the regular point needed in the final tangent
argument.

Keep the domain $D$, local block ring $R$, and nonzero finite
$D$-torsion-free module $L$ of Proposition~\ref{prop:family}.  Thus
$\dim D\geq3$, $D$ is finite over $R$, and
\[
 V=L\widehat\otimes_{\mathcal E}P
\]
is $\OO$-torsion-free and finite over $\mathscr A=\OO[[K]]$.
Here $K$ is a sufficiently deep uniform subgroup of $\SL_2(\Z_2)$.
All coefficient fields below are finite extensions of $\Q_2$.

\subsection{Dimension forces positive rank}

Write $\mathfrak a=(\varpi,\,k-1:k\in K)$ for the maximal
augmentation ideal of $\mathscr A$.  We use its filtration throughout.
After shrinking $K$, the ordered-monomial description of a uniform
Iwasawa algebra gives
\begin{equation}\label{eq:iwasawa-graded}
 \operatorname{gr}_{\mathfrak a}\mathscr A
       \simeq\kk[X_0,X_1,X_2,X_3].
\end{equation}
For $\Z_2$ coefficients this is
\cite[Theorem~3.22 and Lemmas~3.24--3.25]{Venjakob2002}.
The same statement for $\OO$ follows from the ordered expansion in
$\varpi$ and $b_i=k_i-1$, for three ordered uniform generators $k_i$:
take $K$ deep enough that its commutator relations have
$\mathfrak a$-degree at least three.  The symbols of
$\varpi^{a_0}b_1^{a_1}b_2^{a_2}b_3^{a_3}$ then give the polynomial
basis in~\eqref{eq:iwasawa-graded}.

For a nonzero finite $\mathscr A$-module $M$, let $g(M)$ be the degree
of the eventual Hilbert polynomial
\[
 H_M(n)=\operatorname{length}_{\OO}(M/\mathfrak a^nM).
\]
Good filtrations, the Artin--Rees property, and
\eqref{eq:iwasawa-graded} imply $0\leq g(M)\leq4$.
These are the filtered-algebra facts used also in the proof of
\cite[Proposition~4.4]{FM2}.

The decisive comparison is between the dimension of $D$ and this
growth degree. We will make $\dim D-1$ successive prime specializations
to a coefficient fiber chosen to have growth degree at least two:
its reduction modulo $\varpi$ is infinite-dimensional, and $\varpi$
acts injectively. The following lemma shows that each specialization
lowers growth by at least one. Thus $\dim D\geq3$ forces the maximal
growth degree four.

\begin{lemma}\label{lem:growth-drop}
Let $M$ be a finite $\mathscr A$-module and let $f:M\to M$ be
injective and $\mathscr A$-linear.  If $M/fM\ne0$, then
\[
                 g(M/fM)\leq g(M)-1.
\]
Moreover, a finite $\mathscr A$-module of growth degree four has
positive rank over $\mathscr A$.
\end{lemma}
\begin{proof}
The induced filtration on $fM$ is good.  Artin--Rees, followed by the
isomorphism $M\simeq fM$, therefore gives an integer $c\geq1$ such that
\[
 fM\cap\mathfrak a^nM\subseteq f(\mathfrak a^{n-c}M)
 \qquad(n\geq c).
\]
Taking lengths in the induced exact sequence yields
\[
 H_{M/fM}(n)
 \leq H_M(n)-H_M(n-c)=O(n^{g(M)-1}).
\]
This proves the first assertion, including the impossibility of a
nonzero cokernel when $g(M)=0$.

The algebra $\mathscr A$ is a noetherian domain and has an Ore
division ring of fractions.  If a finite module has rank zero,
each element is annihilated by a nonzero element of $\mathscr A$.
A filtration by its finitely many generators has cyclic torsion
quotients.  Such a quotient is of the form $\mathscr A/I$, where
$I$ contains a nonzero element $a$.  Its associated graded module
is a quotient of the polynomial ring in~\eqref{eq:iwasawa-graded}
by the nonzero initial form of $a$, and hence has growth degree at
most three.  The same bound holds for a finite extension of these
quotients.  Thus growth degree four implies positive rank.
\end{proof}

\begin{proposition}\label{prop:positive-rank}
The module $V$ satisfies
\[
                       \rank_{\mathscr A}V>0.
\]
\end{proposition}
\begin{proof}
Choose an auxiliary coefficient point $y$ of $D$ whose local
semisimple parameter, if reducible, has character ratio outside
$\{1,\eps,\eps^{-1}\}$.  Such points exist.  Indeed, at a reducible
exceptional parameter the fixed determinant and the specified ratio
determine each character up to a quadratic character.  Local class
field theory gives only finitely many quadratic characters of
$\Q_2^\times$.  The exceptional parameters therefore lie in a finite
union of coefficient-point closures in $\Spec R$, of integral
dimension at most one.  Finiteness of $D$ over $R$ gives the same
bound for their inverse image.  Since $D$ is horizontal of dimension
at least three, coefficient-point density
\cite[Lemma~3.3]{FM2} supplies $y$ outside that inverse image.

Put $s=\dim D$.  If $\mathfrak p_y$ is the kernel of the coefficient
point, then $D/\mathfrak p_y$ is finite over $\OO$ and has dimension
one.  The complete local domain $D$ is catenary, so there is a
saturated chain
\[
 (0)=\mathfrak p_0\subsetneq\mathfrak p_1
       \subsetneq\cdots\subsetneq\mathfrak p_{s-1}=\mathfrak p_y.
\]
Starting with $L^{(0)}=L$, define $L^{(i+1)}$ by tensoring $L^{(i)}$ with
$D/\mathfrak p_{i+1}$ and removing its torsion over that domain.
Each $L^{(i)}$ is nonzero: the preceding torsion-free finite module has
full support, and Nakayama's lemma at the next prime shows that its
fiber there is nonzero.  Thus $L^{(i+1)}$ has nonzero generic fiber.
The commuting $\mathcal E$-action survives all these operations.

Choose $a_i\in\mathfrak p_{i+1}\setminus\mathfrak p_i$.
Multiplication by $a_i$ is injective on $L^{(i)}$.  The exact compact
block equivalence of Theorem~\ref{input:local-blocks} makes it injective
on
\[
                  V_i=L^{(i)}\widehat\otimes_{\mathcal E}P,
\]
and $V_{i+1}$ is a quotient of its cokernel.  All these objects are
finite over $\mathscr A$, being successive quotients of $V$.
Lemma~\ref{lem:growth-drop} gives
\begin{equation}\label{eq:chain-growth}
             g(V)\geq s-1+g(V_{s-1}).
\end{equation}

Finally, $L^{(s-1)}$ is nonzero and finite torsion-free over $\OO$.
Its associated Banach representation has finite length and, after
a finite coefficient extension, every simple factor has the
central parameter of a coefficient conjugate of $y$; see
Theorem~\ref{input:local-blocks}.
They are infinite-dimensional: the generic principal-series
possibilities are infinite-dimensional, as is the representation
attached to an absolutely irreducible two-dimensional Galois
parameter.  Our choice of $y$ excludes the exceptional character
possibilities.  Hence $V_{s-1}/\varpi V_{s-1}$ is infinite-dimensional
over $\kk$.  Otherwise compact Nakayama would make $V_{s-1}$ finite
over $\OO$, contradicting this Banach description.

Exactness also makes multiplication by $\varpi$ injective on
$V_{s-1}$.  The infinite-dimensional reduction has growth degree at
least one, so Lemma~\ref{lem:growth-drop} gives $g(V_{s-1})\geq2$.
Together with~\eqref{eq:chain-growth}, this proves
\[
                    4\geq g(V)\geq\dim D+1\geq4.
\]
The final assertion of Lemma~\ref{lem:growth-drop} finishes the proof.
\end{proof}

\subsection{Capturing a fixed smooth type}

Let $K_0=\GL_2(\Z_2)$ and let
\[
 \sigma:K_0\longrightarrow\GL_2(\F_2)\simeq S_3
              \xrightarrow{\operatorname{sign}}\{1,-1\}
\]
be the inflated sign character over $E$. Its restriction to the upper
unipotent subgroup of $\GL_2(\F_2)$ is nontrivial, so it is a cuspidal
representation of that finite group. We use this type for two reasons.
Tensoring it with an algebraic representation excludes every Banach
simple with reducible Galois parameter. In a classical smooth representation it
forces depth-zero supercuspidality, and hence a Weil--Deligne parameter
trivial on wild inertia. The second property will exclude the CM fields
inside $\Q(\zeta_{2^\infty})$ and thereby allow the adjoint Selmer
vanishing used in Proposition~\ref{prop:tangent-bound}.

Write
\[
          \Pi_V=\Hom_{\OO}^{\cts}(V,\OO)[1/2]
\]
for the unitary admissible Banach representation associated with $V$.
Its central character on scalar units is $z\mapsto z^{d}$, where
$d=w+1$ by the determinant normalization.  Negative determinant
powers are permitted in the algebraic representations below.

\begin{lemma}\label{lem:type-capture}
There is an irreducible algebraic representation
$U=\Sym^n(E^2)\otimes\det^a$, with $n\geq0$, $a\in\Z$, and
$n+2a=d$, such that
\[
                  \Hom_{K_0}(\sigma\otimes U,\Pi_V)\ne0.
\]
This Hom space is finite-dimensional and carries a continuous
commuting action of $D$.
\end{lemma}
\begin{proof}
Positive rank supplies a nonzero $\mathscr A$-linear map
$V\to\mathscr A$: take a nonzero linear functional after passing
to the Ore division ring and clear right denominators on finitely
many generators.  Maps between finite $\mathscr A$-modules are
continuous.  Choose $K$ so that its intersection with the scalar
unit subgroup $Z_0$ is trivial.  Extend the action on the target
$\mathscr A$ to $Z_0K$ by the dual central character $z^{-d}$.
Since $Z_0K$ has finite index in $K_0$, Frobenius reciprocity gives
a nonzero map from $V$ to its coinduced compact module.  Dualizing
gives a nonzero continuous $K_0$-map
\begin{equation}\label{eq:functions-to-banach}
 \mathcal C_d\longrightarrow\Pi_V,
 \qquad
 \mathcal C_d=\{f\in C(K_0,E):f(kz)=z^{-d}f(k)\}.
\end{equation}
The action on $\mathcal C_d$ is left translation.

We verify the density needed to use this map.  Matrix coefficients
of algebraic representations of central degree zero form an
algebra separating the points of $K_0/Z_0$: coefficients of the
adjoint representation already separate them.  Polynomial
approximation on compact subsets of a finite-dimensional
$2$-adic affine space therefore makes this algebra dense in
$C(K_0/Z_0,E)$.  Concretely, embed the compact quotient by the
adjoint matrix entries and apply multivariable polynomial
approximation on a containing compact box, after extending
continuous functions by clopen partitions.

Fix one algebraic representation of central degree $d$, for example
$\det^{d/2}$ if $d$ is even and $E^2\otimes\det^{(d-1)/2}$ if $d$
is odd.  Its matrix coefficients have no common zero on $K_0$.
On a finite clopen cover of $K_0/Z_0$, divide a function in
$\mathcal C_d$ by a nonvanishing coefficient of this representation
evaluated at $k^{-1}$.  Approximate the resulting functions of central
degree zero as above, and multiply back.  Complete reducibility of
algebraic representations shows that the resulting dense span consists
of matrix coefficients of irreducible algebraic representations with
central degree $d$.  Multiplication by the nowhere-zero function
$\sigma(k^{-1})$ preserves $\mathcal C_d$ and its topology.  Thus
the matrix coefficients of the types $\sigma\otimes U$ in the
statement also have dense span.  A nonzero continuous map in
\eqref{eq:functions-to-banach} cannot kill them all.

For finiteness, transpose a type map to a continuous $K$-equivariant
map from $V$ to the dual of that finite-dimensional type.  A finite
set of $\mathscr A$-generators of $V$ determines such a map, so the
space of maps is finite-dimensional.  Evaluation on those generators
also proves continuity of its commuting $D$-action.
\end{proof}

\begin{lemma}\label{lem:type-irreducible}
Let $E'/E$ be a finite coefficient extension, let $\Pi$ be one of the
$E'$-Banach simples described in Theorem~\ref{input:local-blocks}, and
let $U=\Sym^n(E'^2)\otimes\det^a$, with $n\geq0$ and $a\in\Z$.
If
\[
                  \Hom_{K_0}(\sigma\otimes U,\Pi)\ne0,
\]
then the Galois parameter of $\Pi$ is absolutely irreducible.
\end{lemma}
\begin{proof}
Suppose the parameter is reducible. Theorem~\ref{input:local-blocks}
lists its possible Banach simples. Put
$S_0=\SL_2(\Z_2)$ and
$u=\left(\begin{smallmatrix}1&1\\0&1\end{smallmatrix}\right)$.
On $\sigma\otimes U$, the operator $u$ is minus a unipotent
operator.  In particular $u-1$ is invertible over $E'$.
The group $S_0$ acts transitively on $\mathbb P^1(\Q_2)$.
An $S_0$-map from the type $\sigma\otimes U$ to a continuous principal series, followed
by evaluation at the identity coset, would give a functional invariant
under the upper unipotent subgroup, hence under $u$.  This functional
is zero, and transitivity makes the whole map zero.  The same argument
excludes a character, whose restriction to $S_0$ is trivial.

For a twist of the continuous Steinberg representation, pull back
the defining quotient of a continuous induction by a character along
a proposed type map.  This gives an extension of our finite-dimensional
type by a character.  Every finite-dimensional continuous
representation of $S_0$ over $E'$ is semisimple.  To see this,
continuity and logarithm charts make it analytic on a sufficiently
small open subgroup.  Complete reducibility for
$\mathfrak{sl}_2(E')$ supplies an equivariant projection onto any
invariant subspace.  That projection commutes with a sufficiently
small open normal subgroup of $S_0$; averaging its conjugates over
the finite quotient gives an $S_0$-equivariant projection.
The pulled-back extension therefore splits over $S_0$, producing
the already excluded map into continuous induction.

This excludes every reducible parameter in the classification.
\end{proof}

\subsection{From a type eigenvector to a classical point}

\begin{proposition}\label{prop:regular-point}
There is a coefficient point $x:D\to\OO'$ at which the global
pseudorepresentation is that of a continuous absolutely irreducible
representation
\[
                    r_x:G_{\Q}\longrightarrow\GL_2(E'),
                    \qquad E'=\Frac(\OO'),
\]
after a finite coefficient extension if necessary.  This representation
is odd and finitely ramified.  Its restriction to $G_{\Q_2}$ is
absolutely irreducible and regular de Rham, and
$\WD(r_x|_{G_{\Q_2}})$ is trivial on wild inertia.  Consequently
$r_x$ is attached to a classical cuspidal eigenform up to Tate twist.
\end{proposition}
\begin{proof}
Take a type supplied by Lemma~\ref{lem:type-capture}.  The commuting
$D$-action on its nonzero finite-dimensional Hom space has a common
eigenvector after finite coefficient extension.  It gives a continuous
character $x:D\to E'$.  Since $D$ is compact, its image in the
endomorphism algebra preserves a lattice; its eigenvalues are integral.
Thus $x$ takes values in $\OO'$, and the corresponding map
\begin{equation}\label{eq:type-eigenspace}
                   \sigma\otimes U\longrightarrow\Pi_V[x]
\end{equation}
is nonzero.

This eigenspace has finite length as a Banach representation.
Indeed, extend coefficients to $\OO'$, impose the scalar relations
$a=x(a)$ for $a\in D$ on $L$, and remove $\OO'$-torsion.  The
result is finite over $\OO'$.  The exact compact equivalence identifies
its inverse with the corresponding torsion-free scalar quotient of
$V$; continuous maps to $\OO'$ kill the discarded torsion.
Dualizing and inverting $2$ therefore identifies its Banach
representation with $\Pi_V[x]$.  The finite-length assertion is the
coefficient-fiber assertion of Theorem~\ref{input:local-blocks}.
A composition series now shows that some Banach simple contains
the type in~\eqref{eq:type-eigenspace}: follow a nonzero map through
the series, passing to a quotient when its composite is nonzero and
otherwise factoring it through the closed subrepresentation.
No exactness assertion for the functor of locally algebraic vectors
is involved.

Lemma~\ref{lem:type-irreducible} now makes the local block parameter
$q_x$ absolutely irreducible.  The global determinant at $x$ is then absolutely
irreducible as well and is represented, after finite coefficient
extension, by a continuous $r_x$; this is the representation theorem
for characteristic-zero determinants
\cite[Theorem~2.12]{Chenevier2014}.
Continuity can also be checked by recovering matrix coordinates
from traces against a matrix-algebra basis.
The normalization is
\begin{equation}\label{eq:regular-local-normalization}
                   q_x=r_x|_{G_{\Q_2}}\otimes\eps.
\end{equation}
Since $2$ splits completely in $F$, restriction to $G_F$ is still
absolutely irreducible.

We can now transport the type to the classical forms upstairs.
The point $x$ induces a Hecke coefficient point through
$\mathbb T\to D$.  By \cite[Proposition~4.2]{FM2}, its completed
definite-forms eigenspace contains a Banach tensor product
\[
                         \widehat\bigotimes_{v\mid2}\Pi_v,
\]
where $\Pi_v$ is a simple with parameter $q_x$.  All these parameters
are identical, and an absolutely irreducible parameter has a unique
Banach simple.  Thus each $\Pi_v$ contains the type already found,
and their tensor contains
\[
                         \bigotimes_{v\mid2}(\sigma\otimes U).
\]
These are locally algebraic vectors in the eigenspace of completed
forms.  The classical comparison
\cite[Proposition~4.1]{FM2} identifies them with definite automorphic
forms of algebraic factor $\bigotimes U$.  The weights are regular:
the factor $\Sym^n$ gives classical weight $n+2\geq2$, while the
determinant power only changes the common twist.  In the classical
decomposition the smooth factor contains $\sigma$ at every dyadic
place.  Indeed the Lie algebra identifies the algebraic factor,
and the remaining $K_0$-action is its smooth multiplicity space.
Reduced-norm characters cannot contain this type.  Jacquet--Langlands
therefore gives a regular cuspidal Hilbert eigensystem over $F$.
Its Galois representation agrees with $r_x|_{G_F}$ by the Hecke
identities and Chebotarev.

The finite-group cuspidality of $\sigma$, established above, now
identifies the local smooth factors. The depth-zero type theorem
identifies its characteristic-zero smooth occurrences with depth-zero
supercuspidal representations
\cite[Appendix~A, Sections~A.3.1--A.3.2]{Henniart2002}.
The local Langlands correspondence preserves depth
\cite[Theorem~2.9]{ABPS2016}; its corresponding Weil parameter
therefore has depth zero, which means that it is trivial on wild
inertia.
The dyadic local--global comparison in
\cite[Proposition~4.1]{FM2} now gives regular de Rhamness and this
wild-inertia assertion for $r_x|_{G_F}$ at every dyadic place.
The splitting of $2$ in $F$ gives exactly those properties for
$r_x|_{G_{\Q_2}}$.  The cyclotomic twist in
\eqref{eq:regular-local-normalization} has unramified Weil--Deligne
character and does not change the wild-inertia assertion.

Finally, finite ramification is built into $D$, and its fixed
determinant $\chi$ is odd.  All the hypotheses of
\cite[Theorem~1.1]{FM2} have now been verified for $r_x$: continuity,
absolute irreducibility, oddness, finite ramification, and two
distinct de Rham weights at $2$.  That theorem gives the asserted
classical modularity up to Tate twist, with no restriction on the
residual representation.
\end{proof}

\section{The full Hecke algebra, level, and twists}
\label{sec:hecke}

The regular point produced in Proposition~\ref{prop:regular-point}
is classical up to Tate twist, and its classical level may be divisible
by~$2$.  We now show that both features are compatible with the precise
Hecke algebra of Theorem~\ref{thm:main}.  The same argument will justify
the continuous cyclotomic twist used to normalize the original family.

For this section write $A(N)=\mathbb T_2(N)$ and
$A_k(N)=\Z_2\otimes_{\Z}T_{\leq k}^{(2)}(N)$, as in
\eqref{eq:full-hecke}. Thus an integral polynomial in the Hecke
generators tends to zero if it tends to zero on every fixed finite
range of positive weights.

We use the following structural facts at their stated scope.
The algebra $A(N)$ is reduced and is a finite product of complete
Noetherian local rings with finite residue fields; the topology just
defined induces the maximal-ideal topology on every factor.
It carries a continuous two-dimensional determinant of
$G_{\Q,S}$, where $S$ consists of the primes dividing $2N$, with
\begin{equation}\label{eq:hecke-frobenius}
 t(\operatorname{Frob}_\ell)=T_\ell,
 \qquad d(\operatorname{Frob}_\ell)=\ell S_\ell
 \quad (\ell\nmid 2N).
\end{equation}
Its classical coefficient points are Zariski dense, and the displayed
Hecke elements topologically generate it.  These are
\cite[Lemmas~2.1--2.2 and Proposition~2.3]{Dim}; the Frobenius
convention is exactly that of Equation~(2.2) there.  In particular, the
characteristic polynomial in this convention is
$Z^2-T_\ell Z+\ell S_\ell$.  Finally,
\cite[Theorem~1.1 and Proposition~3.1]{Dim} say, respectively, that
every irreducible component has dimension four, that the closure of
bounded-weight classical points has dimension at most one, and that
the Eisenstein closure has dimension at most two.  It follows that
classical cuspidal points of weight at least three are dense on every
component: remove the latter two closed loci and the other components,
and apply classical density in every nonempty remaining open subset.

These facts remain valid after finite extension of integral
coefficients and passage to a residue factor.  Indeed this extension
is finite flat.  After inverting $2$ it is separable, so reducedness
persists; integral torsion-freeness then gives reducedness before
inverting $2$.  Testing in all coefficient embeddings preserves the
stated density.  Integral finite extensions preserve component
dimensions.  We will use these observations when a character or an
eigenpacket requires larger coefficients.

\subsection{Evaluation on the ordinary tower}

\begin{proposition}\label{prop:hecke-level}
Let $N\geq5$ be odd, and let $f$ be a classical cuspidal eigenform
of positive weight and level dividing $2^aN$, for some $a\geq0$.
Its eigenvalues away from $2N$ define a continuous homomorphism
\[
 A(N)\longrightarrow\OO_{E_f}
\]
after choosing a finite $2$-adic coefficient field $E_f$ containing
them.  The associated determinant has the Frobenius polynomials
of~$f$ in \eqref{eq:hecke-frobenius}.
\end{proposition}

\begin{proof}
We evaluate all the forms on one ordinary tower.  Density of the
tame-level forms on that tower will then turn boundedness of the
Hecke operators into the continuity required here.

\emph{The ordinary tower.}  Enlarge an
integer coefficient ring $\OO$ as necessary.  Over the formal ordinary
locus of the compactified tame modular curve of level $\Gamma_1(N)$,
consider the tower of trivializations
\[
 \iota:\widehat{\mathbb G}_m\xrightarrow{\sim}\widehat E
\]
of the formal group of the universal generalized elliptic curve.
The ordinary locus here includes the cusps.  Write $\mathcal I$ for
the ring of integral functions on this tower, completed for the
coefficient-uniformizer topology, and put
$\mathcal B=\mathcal I[1/2]$, with unit ball $\mathcal I$.
The tower is the inverse system of finite \'etale trivialization
torsors of the ordinary connected $2$-divisible group.  One can see
this by Cartier duality: that group's dual is \'etale of height one.
This also proves the assertion at a multiplicative cusp.
In particular, the tower and $\mathcal I$ are flat over $\OO$.
We use the ordinary-tower construction of
\cite[Section~4.2, Lemma~4.2.1 and Theorem~4.2.2]{Katz1973}.
For Drinfeld level structures we use
\cite[Section~3.2 and Lemma~3.5.1]{KatzMazur1985}; the compactified
moduli interpretation, including cusps in bad characteristic, is
\cite[Theorem~1.2.1 and Definition~2.4.1]{Conrad2007}.

A section of the $i$th power of the Hodge line $\omega$ evaluates
to a function on this tower by expressing its pullback using the
standard differential $dT/(1+T)$ on $\widehat{\mathbb G}_m$.

\emph{Density.} We claim that the $\operatorname{Frac}(\OO)$-linear
span of the evaluations of tame-level classical forms of positive
weights is dense in $\mathcal B$.  Only rational density is needed.

Here are the details at the prime two.  Let $H$ be the weight-one
Hasse invariant on the residue curve.  Its evaluation on the tower
is $1$, and the normalized weight-four Eisenstein series $E_4$
reduces to $H^4$; see \cite[Section~2.1]{Katz1973}.
The open set where $E_4$ is nonvanishing is affine, since $\omega$
is ample.  Its formal completion is the ordinary base; denote its
ring of functions by $\mathcal I_0$.
Serre vanishing for a sufficiently high power of $\omega$ lifts
$H^{4m+1}$ to an integral form $G$ of weight $4m+1$.
Consequently
\[
 \eta=G/E_4^m
\]
is a basis of $\omega$ on the formal ordinary base.  Choose a
parameter $z$ on the formal group whose cotangent is $\eta$.
Such a parameter is obtained by lifting this cotangent successively
on the affine base.  Express the universal trivialization as
\[
 \iota^*z=a_1T+a_2T^2+\cdots.
\]
The ring $\mathcal I$ is topologically generated over $\mathcal I_0$
by the coefficients $a_j$ and $a_1^{-1}$: formal group isomorphisms
are represented by their coefficients, the formal group identities,
and invertibility of the first coefficient. Over the ordinary base
modulo $\varpi^h$, the ideals cutting out the $2^n$-torsion are cofinal
with the powers of the parameter ideal: the formal group has height
one and $2$ is nilpotent on this base. Thus an isomorphism of the formal
groups is a compatible system of isomorphisms of their finite flat
$2^n$-torsion groups. Cartier duality identifies this functor with the
inverse limit of the finite \'etale trivialization torsors. The
coefficient presentation and the tower therefore have the same
coordinate ring modulo every $\varpi^h$. After inverting $2$, formal logarithms
express each individual $a_j$ as a polynomial in $a_1$ with
coefficients in $\mathcal I_0[1/2]$.

Let $\mathcal C$ be the closed linear span under consideration.
It is an algebra, since multiplication adds positive weights.
The evaluations of $E_4$ and $G$ are $1$ modulo the uniformizer.
Their powers with exponents $2^n$ tend to $1$, so $1\in\mathcal C$;
their powers with exponents $2^n-1$ then show that their inverses
also belong to $\mathcal C$.
Ratios of equal-weight sections by powers of $E_4$ give all
functions on the affine base, by the section-ring description of
an ample line bundle.  Completion gives
$\mathcal I_0[1/2]\subseteq\mathcal C$.
Since $\eta=G/E_4^m$ evaluates to $a_1$, both $a_1$ and
its inverse belong to $\mathcal C$.  The logarithm identities now give every
$a_j\in\mathcal C$, and topological generation proves
$\mathcal C=\mathcal B$.
No uniform bound on the denominators of all logarithm coefficients
is required: each coefficient is handled separately before taking
the closure.

\emph{Dyadic level and Hecke operators.} We next evaluate the given form on the same tower and check
that its Hecke operators preserve the integral unit ball.
For forms of level $\Gamma_1(2^aN)$, choose a primitive $2^a$th root of unity in the coefficient ring
and transport its multiplicative torsion point by $\iota$.
This supplies the extra Drinfeld $\Gamma_1(2^a)$-structure.
The construction works on generalized elliptic curves at cusps:
the original tame level is ample, and adding a subgroup in the
identity component preserves ampleness.  The formal group of a
Tate polygon is $\widehat{\mathbb G}_m$ by
\cite[Equation~(2.5.3)]{Conrad2007}.  The integral sections of the
Hodge line recover classical forms after inverting $2$
\cite[Section~4.4, Equation~(4.4.2)]{Conrad2007}; hence a fixed
scalar multiple of a classical form extends over the integral
model and has evaluation in $\mathcal I$.  Evaluation is nonzero for a
nonzero form, as its $q$-expansion at a compatible multiplicative
cusp shows.

The good Hecke operators preserve $\mathcal I$ and commute with
all these evaluations.  For $T_\ell$ the correspondence sums over
degree-$\ell$ isogenies, with normalization $\ell^{-1}$, a unit
in $\OO$.  The isogenies induce isomorphisms of formal groups,
and the finite flat correspondence and its trace are integral.
The extension over multiplicative cusps, and the normalization
of trace--pullback as $\ell T_\ell$, are
\cite[Theorem~4.4.3 and Equation~(4.5.1)]{Conrad2007}.
For $\ell S_\ell$, the diamond operator and the scalar change
of trivialization contribute $\ell^i\ell^{-1}$ in weight $i$,
which is the required $\ell^{i-1}\langle\ell\rangle$.
This includes the diamond action on the additional $2$-power
level.  Thus every integral polynomial in these operators is
a contraction of $\mathcal B$.

\emph{Continuity.} Let $h_\alpha$ be a net of such polynomials
that tends to zero at each finite weight stage.  On every finite
sum of evaluated tame classical forms, $h_\alpha$ tends to zero.
Density and the uniform contraction bound imply
\[
 h_\alpha v\longrightarrow0\qquad(v\in\mathcal B):
\]
first approximate $v$ by a finite sum, then use the contraction
bound on the error.  Apply this to the nonzero evaluation $v_f$
of~$f$.  Since $h_\alpha v_f=\lambda_f(h_\alpha)v_f$,
the eigenvalues tend to zero.  The same argument shows that a
polynomial relation in $A(N)$ acts as zero on $v_f$.
The eigenvalue map therefore extends continuously from the dense
polynomial algebra to $A(N)$.  Its values are integral, either
by the contraction bound or by classical integrality of good
Hecke eigenvalues.  Compatibility of evaluation with the Hecke
action proves the asserted Frobenius identities.
\end{proof}

\subsection{Continuous cyclotomic twists}

\begin{proposition}\label{prop:hecke-twists}
Let $N\geq5$ be odd.  Let $\lambda:A(N)\to\OO_E$ be a continuous
coefficient point, with associated two-dimensional determinant
$(t_\lambda,d_\lambda)$.  If $E'/E$ is finite and
$\theta:\Z_2^\times\to\OO_{E'}^\times$ is continuous, then
the twisted determinant
\[
 \bigl((\theta\circ\eps)t_\lambda,
       (\theta\circ\eps)^2d_\lambda\bigr)
\]
is also supplied by a continuous coefficient point of $A(N)$.
In particular this holds for every integer Tate twist, with the
same tame level $N$.
\end{proposition}

\begin{proof}
Fix integral coefficients containing the residual values, and
work in the factor $A$ of the resulting Hecke algebra selected
by $\lambda$.  Let $R^{\mathrm{ps}}$ be the universal global
two-dimensional determinant deformation ring for these residual
data, with ramification allowed at the primes dividing $2N$ and
with varying determinant.  The structural facts above give a
continuous surjection
\[
 R^{\mathrm{ps}}\twoheadrightarrow A,
 \qquad J=\ker(R^{\mathrm{ps}}\longrightarrow A).
\]

First suppose $\theta$ has finite order, and enlarge coefficients
to contain its values.  Twisting the universal determinant over
$A$ defines a continuous map $R^{\mathrm{ps}}\to A$ over these
enlarged coefficients.  There is no change of residual factor:
every finite quotient of $\Z_2^\times$ is a $2$-group, so a
finite-order character reduces to $1$ in characteristic two.
At every classical cuspidal point of weight at least three,
the twisted determinant belongs to a classical finite-character
twist.  Its extra conductor is a power of $2$; Proposition~\ref{prop:hecke-level}
therefore supplies a coefficient point of the same $A(N)$.
For $j\in J$, its image under the twisted map vanishes at all
these points.  Their density and reducedness show that this image
is zero.  Thus twisting factors through $A$ and preserves all its
coefficient points.  Equality of the specialized determinants
follows from \eqref{eq:hecke-frobenius} and Chebotarev.

Now fix $\lambda$, and fix a sign $s\in\{1,-1\}$.  Every
$u\in\Z_2^\times$ has a unique expression $u=(-1)^e5^b$,
with $e\in\{0,1\}$ and $b\in\Z_2$.  Define the universal character
\[
 \Theta_X((-1)^e5^b)=s^e(1+X)^b
 \quad\text{in }\OO_{E'}[[X]]^\times.
\]
Twisting the determinant of $\lambda$ by $\Theta_X\circ\eps$
gives a continuous map
$R^{\mathrm{ps}}\to\OO_{E'}[[X]]$.
For every $2$-power root of unity $\zeta$, evaluation at
$X=\zeta-1$ is a finite-order twist, so the image of every
$j\in J$ vanishes at all such values.  A nonzero series in
$\OO_{E'}[[X]]$ cannot do this.  Indeed, after dividing by the
largest common uniformizer power of its coefficients,
Weierstrass preparation expresses it as a distinguished polynomial
times a unit.  The polynomial has only finitely many roots in
the open unit disk.  Therefore the image of $J$ is zero, and the
universal twisted map factors continuously through $A$.

For the given continuous $\theta$, its image on the pro-$2$
group $1+4\Z_2$ reduces trivially in the odd-order group
$\kk_{E'}^\times$.  Hence $\theta(5)-1$ lies in the maximal
ideal of $\OO_{E'}$.  Specialize $s=\theta(-1)$ and
$X=\theta(5)-1$.  This continuous specialization gives the
required point.  All ring topologies used here are the completed
local topologies, which agree with the prescribed finite-stage
topology by \cite[Proposition~2.3]{Dim}.
\end{proof}

We can now pass from classical representations in the companion
normalization to points of $A(N)$.  If $V_f$ denotes the arithmetic representation of a
weight-$k$ primitive form with nebentype $\nu$, then
$\det V_f=\nu\eps^{k-1}$, and the wedge pairing gives
\begin{equation}\label{eq:classical-dual-conversion}
 V_f^\vee\simeq V_f\otimes\nu^{-1}\eps^{1-k}
       \simeq V_g\otimes\eps^{1-k},
\end{equation}
where $g$ is the primitive form associated with
$f\otimes\nu^{-1}$.  Thus an expression
$r_x\simeq V_f^\vee\eps^n$ becomes
$r_x\simeq V_g\eps^{n+1-k}$.
This is the conversion in \cite[Section~2.1]{FM2}.
Propositions~\ref{prop:hecke-level} and~\ref{prop:hecke-twists}
therefore put every representation classical up to Tate twist
on $A(N)$ for a sufficiently divisible odd $N$.

\section{From the regular point to the original representation}
\label{sec:conclusion}

We have a fixed-determinant trace domain $D$ containing both the
normalized original point and the regular point $x$ of
Proposition~\ref{prop:regular-point}.  The latter point is now on
Hecke support over $\Q$.  To put the whole domain on that support,
we compare both quotients of the global deformation ring in which
the determinant is allowed to vary.  The decisive fact is that this
larger ring is regular at~$x$.

\subsection{The full adjoint tangent space}

\begin{proposition}\label{prop:tangent-bound}
Let $r_x:G_{\Q,S}\to\GL_2(E_x)$ be odd, absolutely irreducible,
and classical cuspidal up to Tate twist.  Suppose its restriction
to $G_{\Q_2}$ is absolutely irreducible and de Rham with distinct
Hodge--Tate weights, and its Weil--Deligne representation is trivial
on wild inertia.  Here $S$ is any finite set containing $2$ and all
ramified primes.  Then
\[
 \dim_{E_x}H^1(G_{\Q,S},\ad r_x)\leq3.
\]
Consequently, if $R_{\mathrm{glob}}$ is the universal global
determinant deformation ring with varying determinant for the
semisimplified residual data of $r_x$, and $\mathfrak x$ is its
coefficient prime, then
\[
 \edim (R_{\mathrm{glob}})_{\mathfrak x}\leq3.
\]
Both assertions permit scalar or reducible semisimplified residual
data and finite extension of the coefficient ring.
\end{proposition}

\begin{proof}
Set $W=\ad r_x=\End_{E_x}(r_x)$, the full four-dimensional
adjoint representation. Tameness of the Weil--Deligne parameter will
give the CM qualification for adjoint Selmer vanishing; local
irreducibility and distinct de Rham weights will give the local
cohomology bound. For a finite place $v$, the finite local
condition is
\[
 H^1_f(\Q_v,W)=
 \begin{cases}
  \ker\bigl(H^1(\Q_v,W)\to H^1(I_v,W)\bigr),&v\ne2,\\
  \ker\bigl(H^1(\Q_2,W)\to
        H^1(\Q_2,W\otimes_{\Q_2}B_{\mathrm{cris}})\bigr),&v=2.
\end{cases}
\]
Here $B_{\mathrm{cris}}$ is Fontaine's crystalline period ring.
The global group $H^1_f(\Q,W)$ imposes these conditions at all
finite places.  Positive-degree cohomology at the real place is
zero over $E_x$.

Newton--Thorne's theorem \cite[Theorem~5.4]{NewtonThorne2023},
specialized to $\GL_2/\Q$, states that for a regular algebraic
cuspidal automorphic representation $\pi$ and any prime $p$,
\[
 H^1_f(\Q,\ad r_{\pi,p})=0
\]
provided either $\pi$ is non-CM, or its CM field $K$ is not
contained in $\Q(\zeta_{p^\infty})$.  This statement uses the
full adjoint and has no residual irreducibility hypothesis.
Our classical representation is regular algebraic, including
when its classical weight is two, and twists do not change its
adjoint.  It remains only to verify the CM qualification.

If its CM field were contained in $\Q(\zeta_{2^\infty})$, the
corresponding quadratic self-twist would restrict to a nontrivial
wildly ramified character at $2$.  Indeed the three quadratic
subfields are $\Q(i)$, $\Q(\sqrt2)$, and $\Q(\sqrt{-2})$;
each ramifies at $2$, and a quadratic ramified character there
is necessarily wild.  The Weil--Deligne parameter would then
be isomorphic to its twist by that character.  On a wild inertia
element where the character is $-1$, these two parameters act
as $I$ and $-I$, respectively, by the hypothesis on~$x$.
This is impossible.  Newton--Thorne therefore gives
\begin{equation}\label{eq:global-adjoint-selmer}
 H^1_f(\Q,W)=0.
\end{equation}

We next compute the local quotients.  The trace pairing identifies
$W^*$ with $W$, so local Tate duality identifies the dual of
$H^2(\Q_v,W)$ with $\Hom_{G_{\Q_v}}(r_x,r_x(1))$.
At $2$, a nonzero such map is an isomorphism by absolute
irreducibility; its determinants would force $\eps^2=1$ on
$G_{\Q_2}$. Thus $H^2(\Q_2,W)=0$, including in weight two.
At $v\ne2$, local--global compatibility
\cite[Theorem~A]{Carayol1986} and genericity of each local factor
of a cuspidal automorphic representation give the same vanishing:
the genericity criterion
\cite[Lemma~1.1.3]{Allen2016} says that its Frobenius-semisimple
Weil--Deligne parameter $U$ satisfies
$\Hom_{\mathrm{WD}}(U,U(1))=0$.
A Galois morphism would induce such a morphism, also after
Frobenius semisimplification.  Tate twists do not affect these
adjoint computations.  This is the non-dyadic argument of
\cite[Lemma~5.2]{Dim}.

Write $h^j_v=\dim_{E_x}H^j(\Q_v,W)$.
For $v\ne2$, local Euler characteristic gives $h^1_v=h^0_v$,
whereas
\[
 H^1_f(\Q_v,W)=W^{I_v}/(\operatorname{Frob}_v-1)W^{I_v}
\]
also has dimension $h^0_v$.  Hence the local quotient is zero.
At $2$, Schur's lemma gives $h^0_2=1$, and the Euler
characteristic formula gives $h^1_2=1+\dim W=5$.
The filtration degrees of the filtered de Rham module
$D_{\mathrm{dR}}(W)$ are
$0,0,h,-h$ for some nonzero integer $h$.  The Bloch--Kato
dimension formula \cite[Corollary~3.8.4]{BlochKato1990} gives
\[
 \dim H^1_f(\Q_2,W)
 =h^0_2+
   \dim D_{\mathrm{dR}}(W)/\operatorname{Fil}^0=1+1=2.
\]
We have proved
\begin{equation}\label{eq:local-adjoint-quotients}
 \dim H^1(\Q_v,W)/H^1_f(\Q_v,W)
 =\begin{cases}0,&v\ne2,\\3,&v=2.\end{cases}
\end{equation}
The kernel of the restriction map
\[
 H^1(G_{\Q,S},W)\longrightarrow
 \bigoplus_{v\in S} H^1(\Q_v,W)/H^1_f(\Q_v,W)
\]
lies in $H^1_f(\Q,W)$, because classes in its source are already
unramified outside $S$.  Equations~\eqref{eq:global-adjoint-selmer}
and~\eqref{eq:local-adjoint-quotients} now give the bound three.

For the embedding-dimension assertion, apply the characteristic-zero
tangent comparison of \cite[Proposition~4.1]{Dim} to
$R_{\mathrm{glob}}$. Its trace and determinant values topologically
generate the ring, and the specialization at $\mathfrak x$ is
absolutely irreducible. The comparison bounds the embedding dimension by
$\dim_{E_x}H^1(G_{\Q,S},W)$, with no condition on the
semisimplified residue. The continuity needed here follows because
$R_{\mathrm{glob}}/\mathfrak x^2$ is finite over $\OO$:
$R_{\mathrm{glob}}/\mathfrak x$ is finite over $\OO$, and
$\mathfrak x/\mathfrak x^2$ is finite over that quotient.
Finite extension of $\OO$ adds no tangent directions, since
characteristic-zero derivations kill its fraction field by
separability over $\Q_2$. The preceding cohomological bound therefore
gives the asserted embedding dimension.
\end{proof}

\subsection{A regular point determines the supporting component}

The following elementary lemma isolates the use of the
four-dimensional Hecke theorem.  It explains why the tangent
calculation is needed in addition to the existence of a modular point.

\begin{lemma}\label{lem:regular-support}
Let $R_0$ be a complete Noetherian local $\OO$-algebra, and let
$A_0$ and $D_0$ be quotients of $R_0$, with $D_0$ a domain.
Suppose every irreducible component of $\Spec A_0$ has dimension
four.  Let $x_0:R_0\to\OO'$ be a continuous coefficient point
factoring through both quotients, where $\OO'/\OO$ is a finite
extension of integer coefficient rings.  If
\[
 \edim(R_0)_{\mathfrak x_0}\leq3,
 \qquad \mathfrak x_0=\ker x_0,
\]
then $R_0\to D_0$ factors through $A_0$.
\end{lemma}

\begin{proof}
Choose a minimal prime $P$ of $A_0$ contained in its prime
$\mathfrak x_0$.  The complete local domain $A_0/P$ is
catenary.  Moreover $A_0/\mathfrak x_0$ is an $\OO$-submodule
of $\OO'$, hence finite over $\OO$, and it contains $\OO$.
It therefore has dimension one.  The dimension formula gives
\[
 \dim(A_0/P)_{\mathfrak x_0}
 =\dim A_0/P-\dim A_0/\mathfrak x_0=4-1=3.
\]
As $R_0$ surjects onto $A_0$, we obtain
\[
 3\leq\dim (A_0)_{\mathfrak x_0}
 \leq\dim (R_0)_{\mathfrak x_0}
 \leq\edim (R_0)_{\mathfrak x_0}\leq3.
\]
Thus $(R_0)_{\mathfrak x_0}$ is a regular local ring of dimension
three, in particular a domain.  Its quotient
$(A_0)_{\mathfrak x_0}$ has the same dimension, so the localized
kernel $\ker(R_0\to A_0)_{\mathfrak x_0}$ is zero: a nonzero
ideal in this regular local domain has quotient of smaller
dimension.

If $a\in\ker(R_0\to A_0)$, there is consequently
$s\notin\mathfrak x_0$ with $sa=0$ in $R_0$.
The kernel of $R_0\to D_0$ is contained in $\mathfrak x_0$,
so the image of $s$ in $D_0$ is nonzero.  Since $D_0$ is a
domain, the image of $a$ must vanish.  This proves the factorization.
\end{proof}

\begin{proof}[Proof of Theorem~\ref{thm:main}]
Apply Proposition~\ref{prop:family} to the original representation,
after its indicated continuous cyclotomic normalization and finite
coefficient extension.  It supplies the horizontal trace domain
$D$ with the normalized original coefficient point.
Propositions~\ref{prop:diagonal}, \ref{prop:positive-rank},
and~\ref{prop:regular-point} supply a coefficient point $x$ of
$D$ satisfying all the hypotheses of
Proposition~\ref{prop:tangent-bound}.

Choose an odd $N\geq5$ sufficiently divisible that it contains
every odd ramified prime of the original family and that the
classical representation underlying $r_x$ has level dividing
$2^aN$ for some $a$.  Increase the allowed global ramification
set to the primes dividing $2N$; the existing family remains
unramified at any newly added primes.  By
Propositions~\ref{prop:hecke-level} and~\ref{prop:hecke-twists},
$r_x$ gives a continuous point of $\mathbb T_2(N)$ with the
required Frobenius values.

Fix the coefficient ring $\OO$ of $D$ and the residual embedding
selected by~$x$.  Let $A$ be the corresponding residue factor
of $\OO\otimes_{\Z_2}\mathbb T_2(N)$, and let
$R_{\mathrm{glob}}$ be the universal global determinant deformation
ring for this residual determinant with \emph{varying} determinant.
Let $\OO_x$ be a coefficient ring for $x$. The determinant on $D$
and the Hecke determinant give a commutative diagram whose top
arrows are continuous surjections:
\[
 \begin{array}{ccc}
   &R_{\mathrm{glob}}&\\[-2pt]
   \swarrow&&\searrow\\[-2pt]
   D&&A\\[-2pt]
   \searrow&&\swarrow\\[-2pt]
   &\OO_x.&
 \end{array}
\]
The lower arrows are the coefficient point $x$; equality follows
from its Frobenius identities and Chebotarev.  The quotient $D$
has fixed determinant $\chi$, whereas $R_{\mathrm{glob}}$ and
$A$ retain determinant variation.  Surjectivity onto $A$ follows
from topological generation by the Hecke elements and compactness,
as in \cite[Proposition~2.3]{Dim}; surjectivity onto $D$ is its
definition as a trace-image quotient.

Every component of $A$ has dimension four by
\cite[Theorem~1.1]{Dim} and finite coefficient extension.
Proposition~\ref{prop:tangent-bound} gives
$\edim(R_{\mathrm{glob}})_{\mathfrak x}\leq3$.
Lemma~\ref{lem:regular-support} therefore makes
$R_{\mathrm{glob}}\to D$ factor through $A$.
The normalized original coefficient point of $D$ is consequently
a point of the required Hecke algebra over~$\Q$.

Undo the normalization with Proposition~\ref{prop:hecke-twists}.
We obtain, over a finite extension $E'/E$, a continuous map
\[
 \lambda:\mathbb T_2(N)\longrightarrow\OO_{E'}
\]
whose values on $T_\ell$ and $\ell S_\ell$ are respectively
$\tr r(\operatorname{Frob}_\ell)$ and
$\det r(\operatorname{Frob}_\ell)$ for $\ell\nmid2N$.
The compact image of $r$ preserves an $\OO_E$-lattice, so all
these values lie in $\OO_E$.  Since the Hecke generators
topologically generate $\mathbb T_2(N)$ and $\OO_E$ is closed
in $\OO_{E'}$, the entire image of $\lambda$ lies in $\OO_E$.
The induced map to $\OO_E$ is continuous for its $2$-adic topology.
Our choice of $N$ includes every odd ramified prime of the
original $r$, which completes the proof.
\end{proof}


\bibliographystyle{plain}
\bibliography{refs}
\end{document}
